% P34 fresh Korean production translation.
% Initial translation authority: ED0005 / NOETH-DE-AUTH-v038-20260805.
% Current scope-rebase authority: ED0011 / NOETH-DE-AUTH-v047-20260808.
% Scope: current German lines 16335--18360.
% Complete translated scope is integrated in reader.tex.

\providecommand{\mE}{\mathfrak E}
\providecommand{\mA}{\mathfrak A}
\providecommand{\mB}{\mathfrak B}
\providecommand{\mC}{\mathfrak C}
\providecommand{\mD}{\mathfrak D}
\providecommand{\mG}{\mathfrak G}
\providecommand{\mH}{\mathfrak H}
\providecommand{\mK}{\mathfrak K}
\providecommand{\mL}{\mathfrak L}
\providecommand{\mM}{\mathfrak M}
\providecommand{\mN}{\mathfrak N}
\providecommand{\mO}{\mathfrak O}
\providecommand{\mP}{\mathfrak P}
\providecommand{\mR}{\mathfrak R}
\providecommand{\mS}{\mathfrak S}
\providecommand{\mT}{\mathfrak T}
\providecommand{\mU}{\mathfrak U}
\providecommand{\mZ}{\mathfrak Z}
\providecommand{\mo}{\mathfrak o}
\providecommand{\ma}{\mathfrak a}
\providecommand{\mb}{\mathfrak b}
\providecommand{\mc}{\mathfrak c}
\providecommand{\md}{\mathfrak d}
\providecommand{\mf}{\mathfrak f}
\providecommand{\mh}{\mathfrak h}
\providecommand{\ml}{\mathfrak l}
\providecommand{\mn}{\mathfrak n}
\providecommand{\mr}{\mathfrak r}
\providecommand{\mt}{\mathfrak t}
\providecommand{\mx}{\mathfrak x}
\providecommand{\mz}{\mathfrak z}
\providecommand{\iso}{\simeq}
\providecommand{\ann}{\operatorname{Ann}}
\providecommand{\tuple}[1]{(#1)}
\providecommand{\mat}[1]{\begin{pmatrix}#1\end{pmatrix}}
\providecommand{\glqq}{“}
\providecommand{\grqq}{”}

\section*{34. 초복소량과 표현론}
\emph{Math. Zs. 30 (1929), 641--692쪽.}

\begin{center}
{\Large\bfseries 초복소량과 표현론\footnote{이하는 B. L. van der Waerden이 나의 1927/28년 겨울학기 강의를 자유롭게 정리한 것이다. 출판을 위한 원고는 우리가 함께 다듬었다. 또한 여러 비판적 논평을 해 준 B. L. van der Waerden에게 감사를 표한다.}.}\\[1em]
글쓴이\\[0.5em]
괴팅겐의 에미 뇌터.\\[0.8em]
\rule{3em}{0.4pt}
\end{center}

\subsection*{서론.}

초복소계에 관한 가장 중요한 일반정리들은 Molien(Math. Annalen 41권 및 Dorpater Ber. 1897)에게서 비롯한다. 그와 본질적으로 독립하여 프로베니우스는 곧이어 초복소계와 그 표현의 이론, 특히 유한군의 표현론을 통일적으로 전개하였다. 그 기초는 데데킨트에게서 유래한 군행렬식, 더 일반적으로는 임의의 초복소계의 행렬식이라는 개념이다. 프로베니우스는 이 행렬식의 서로 다른 기약인자들에 --- 계수영역은 언제나 복소수 전체의 체이다 --- 서로 다른 기약표현류들이 대응하며, 이 방식으로 모든 기약표현류가 빠짐없이 얻어진다는 것을 보였다. 그러한 표현류는 그 \glqq{}지표계\grqq{}로 완전히 특징지어진다. 유한군의 경우 이 지표계는 군의 켤레원소류들에서 생성되는 가환 초복소계의 행렬식을 분해하여 얻어진다. 이 행렬식은 지표들을 계수로 하는 일차인자들로 분해된다. 이는 유한 아벨군의 군행렬식이 그 아벨군의 서로 다른 지표들을 계수로 하는 일차인자들로 분해된다는, 데데킨트가 처음 얻은 결과의 직접적인 일반화이다(프로베니우스와의 서신교환).

그러나 프로베니우스는 개념적으로 단순하고 투명한 이 결과들을 힘겨운 계산으로 얻었다. 후대의 발전은 이 결과들을 더 단순하게 유도하는 동시에 임의의 체를 계수영역으로 삼는 경우로 확장하는 데 향하였다. 이 발전은 초복소계와 표현론에서 완전히 따로 진행되었다. 초복소계는 웨더번에게서 산술적으로 다루어졌다. 모든 초복소계는 양쪽으로 직접분해되지 않는 환들의 직접합으로 유일하게 분해된다. 근기가 없는 계에서는 이 직접분해되지 않는 환들이 본질적으로 유일하게 결정되는 한 대응체 --- 반드시 가환일 필요는 없다 --- 의 원소로 된 모든 \(n\)행 정방행렬의 계와 동형이다\footnote{Dickson, \emph{Algebras and their Arithmetics}(독일어판 \emph{Algebren und ihre Zahlentheorie}, Zürich 1927)의 문헌목록을 보라.}. 표현론은 Burnside와 I. Schur에게서 초복소계와 무관하게 주어진 표현에서 직접 출발하여 행렬정리들을 쓰는 초등적 기초를 얻었다. 특히 Schur는 주어진 한 체에서 기약인 표현이 절대적으로 분해되는 최소차수 수체를 물었다. 절대기약 성분들은 유한개의 켤레 표현류에 속한다. 찾는 수체들의 차수는 이 표현류의 수와 인덱스, 곧 각 켤레류에 속하는 성분의 수의 곱과 같다. Schur 이론의 주된 도구는 한 기약표현과 가환하는 모든 행렬이 이루는 계이다\footnote{Speiser의 군론 제10장의 문헌목록을 보라.}.

이하의 순수하게 산술적인 정초에서는 초복소량과 표현론이 다시 하나의 통일된 전체로 나타난다. 이는 일정한 유한성 조건만을 만족하는 비가환환의 일반이론의 특수한 경우이다. 정확히는 이 환들에 관한 가군류와 아이디얼류의 이론이며, 주된 결과는 기약 가군류들이 이미 대응 아이디얼류들로 모두 얻어진다는 것이다. 특히 근기가 없는 환에서는 모든 가군류가 기약인 것들로 분해되어 완전가약이 된다. 이는 정칙 군행렬식 또는 계행렬식의 기약인자들만으로 이미 모든 기약표현류를 얻고, 근기가 없는 계에서는 표현들이 완전가약이라는 프로베니우스 결과의 산술적 동치이다. 실제로 계행렬식과 그 분해를 고찰하는 것은 노름으로 넘어가는 일로 볼 수 있지만, 여기서는 아이디얼의 직접합분해와 그 합성열을 직접 다룬다. 표현론은 표현가군이라는 개념을 통해 가군류의 이론이 된다.

가군류와 아이디얼류의 도입은 순전히 군론적으로 정초된다. 가군과 아이디얼은 덧셈에 관해서는 아벨군이되, 환원소와의 일정한 곱셈을 허용한다는 제한을 받는 것으로 이해한다. 즉 이들은 \glqq{}연산자를 갖는 군\grqq{}을 이룬다(§ 1). 이러한 군에서는 통상적인 동형 대신 \glqq{}연산자동형\grqq{}이 나타난다(§ 2). 한 고정된 영역 안에서 서로 연산자동형인 군들을 한 류로 묶는다. 이 류 개념은 예컨대 수체의 아이디얼의 경우 통상적인 류 개념과 일치한다. 연산자를 갖는 군의 이론은 W. Krull과 O. Schmidt에게서 유래하며(각주 6 참조), 제I장에서 체계적으로 전개한다. 합성열, 직접곱 또는 직접합, 완전가약군에 관한 정리들도 여기서 여전히 성립한다. 위의 류 구분이라는 의미에서의 이 유일성정리들이 이후 모든 논의의 기초이다.

우선 이 정리들로부터 \emph{임의의 표현}의 기약 대각성분에 관한 일반 \emph{유일성정리}를 얻는다(제III장 § 16). 이는 표현류와 표현가군류, 곧 연산자를 갖는 군의 류가 일대일로 대응한다는 사실에 기초한다. 표현들의 \emph{전체}에 관한 더 나아간 물음은 표현되는 환, 특수하게는 초복소계의 구조를 더 자세히 살필 것을 요구한다. 제II장에서는 한쪽아이디얼의 고찰을 전면에 두는 \emph{군론적} 관점에 기초하여 웨더번의 결과들을 새로 얻고 발전시킨다. 구체적으로 우측아이디얼에 대한 \glqq{}배수사슬정리\grqq{}, 또는 그와 같은 \glqq{}최소조건\grqq{} --- 우측아이디얼들의 모든 집합에는 그 집합 안에서 극소인 것이 적어도 하나 있다 --- 만으로 유한성 조건이 충분하다는 것이 드러난다\footnote{\glqq{}이중사슬정리\grqq{}를 가정하면 웨더번의 추론을 옮길 수 있다. E. Artin, \emph{Zur Theorie der hyperkomplexen Zahlen}, Hamb. Abh. 1927을 보라. 또 A. Suschkewitsch, \emph{Über die endlichen Gruppen ohne das Gesetz der eindeutigen Umkehrbarkeit}, Math. Annalen 99 (1928), 30--50쪽을 보라. 여기서는 역산할 수 없는 단 하나의 결합적 연산만을 갖는 유한 영역에 대해 나란한 구조정리들을 전개한다. Suschkewitsch가 \glqq{}핵\grqq{}을 우측군과 좌측군으로 분해하는 것은 \emph{단위원을 갖는 양쪽단순환}을 우측아이디얼과 좌측아이디얼들의 합으로 나타내는 것에 대응한다.}. 근기가 없고 최소조건을 만족하는 환들은 단위원을 갖는 우측 완전가약환들과 정확히 같다. 그러한 환에서 양쪽으로 직접분해되지 않는 한 환의 모든 단순 우측아이디얼은 같은 류에 속한다. 따라서 연산자를 갖는 군으로서 같은 자기동형환을 가지며, 아이디얼이 단순하므로 이 환은 체가 된다. 바로 이 아이디얼류의 자기동형체가 웨더번의 행렬표현을 매개한다.

자기동형체에서의 이 행렬표현은 완전가약인 경우에 자기동형체 또는 그 부분체들에 관한 표현을 요구하면 표현문제를 동시에 해결한다. 특히 초복소계의 계수영역에 관한 경우가 그렇다. 가군류를 아이디얼류로 환원하는 정리의 직접적인 결과로, 웨더번의 표현과 그 자기동형체의 원소들을 자연스러운 방식으로 한 부분체의 행렬들로 바꾸어 얻는 표현 이외에는 다른 기약표현이 없다.

따라서 완전가약인 경우 서로 다른 기약표현류의 수는 아이디얼류의 수와 같고, 이는 서로 다른 양쪽 직접분해불능환, 곧 단순환의 수와 같다. 마지막 수는 다시 중심의 직접분해불능 성분들의 수와 같은데, 이 성분들은 가환체가 된다. 대수적으로 닫힌 계수영역 위의 초복소계에서는 이 가환체들이 중심의 서로 다른 환준동형, 곧 기약표현들로 완전히 결정된다. 이것들은 한 \emph{수인자}를 제외하면 지표들이다.

얻은 결과들은 근기를 갖는 계의 기약표현도 개관하게 해 준다. 이러한 표현은 근기에 따른 잉여류환의 표현으로 환원되며, 그 잉여류환은 근기가 없는 계가 되기 때문이다.

\emph{뒤에서 보이듯이}, 근기가 없는 초복소계에서는 자기동형체로부터 \glqq{}분해체\grqq{}의 이론도 나온다. 분해체란 계수영역의 가환확대 가운데 표현들이 절대기약인 것들로 분해되는 것, 달리 말해 절대단순 우측아이디얼들의 직접합분해가 일어나는 것이다. \emph{분해체}들은 자기동형체의 분해체들과 같고, \emph{최소차수의 모든 분해체}는 자기동형체의 \emph{극대 가환부분체}와 동형이다. 위에서 말한 Schur의 연구와의 관계는 주어진 표현과 가환하는 행렬들의 \emph{전치행렬}들이 바로 자기동형체의 한 표현을 준다는 사실로 세워진다\footnote{R. Brauer--E. Noether, \emph{Über minimale Zerfällungskörper irreduzibler Darstellungen}, Sitz.-Ber. Preuß. Akad. Wiss. 1927, 221쪽을 보라.}. 이 정리들은 비가환체의 갈루아 이론이라는 틀 안에서 체계적으로 전개할 예정이다.

바탕에 둔 개념들, 곧 연산자준동형과 자기동형환은 최소조건만을 만족하는 근기 있는 일반환의 구조도 낳는다. 이는 다른 곳에서 전개될 것이다.

\section*{제I장. 군론적 기초.}

\subsection*{§ 1. 연산자를 갖는 군\footnotemark.}
\footnotetext{여기서 설명하는 개념들은 W. Krull, \emph{Über verallgemeinerte endliche Abelsche Gruppen}, Math. Zeitschr. 23 (1925), 161--196쪽과 O. Schmidt, \emph{Über unendliche Gruppen mit endlicher Kette}, Math. Zeitschr. 29 (1928), 34--41쪽에서 유래한다.}

\(\mG\)를 유한 또는 무한 군이라 하고 원소를 \(a,b,\ldots\)라 쓰자. 군 \(\mG\)의 연산자영역 \(\Omega\)란 새 기호 \(H,\Theta,\ldots\)의 집합으로서, 모든 \(a\in\mG\), \(\Theta\in\Omega\)에 대해 \(\Theta a\in\mG\)가 유일하게 정의되고 분배법칙
\[
        \Theta(ab)=\Theta a\cdot\Theta b
\]
이 성립하는 것이다. 따라서 각 연산자는 군에서 자기 자신으로 가는 준동형을 정의하며, 군은 자기 자신 또는 한 부분군으로 사상된다. 단위원은 단위원으로, 역원은 역원으로 간다.

서로 다른 연산자들이 서로 다른 준동형을 정의하면 그 연산자영역을 \emph{절대적}이라 한다. 임의의 연산자영역에서 같은 준동형을 생성하는 모든 원소를 같게 놓으면 절대 연산자영역을 얻는다. 절대 연산자영역은 군에서 자기 자신으로 가는 모든 준동형의 집합의 한 부분집합과 일대일로 대응한다.

고정된 연산자영역 \(\Omega\)에 관하여, 모든 \(a\in\mH\), \(\Theta\in\Omega\)에 대해 \(\Theta a\in\mH\)인 부분군 \(\mH\)를 \(\mG\)의 \emph{허용부분군}이라 한다.

\emph{예.} 1. 연산자들이 내부자기동형 \(\Theta a=c^{-1}ac\)라 하자. 허용부분군은 정규부분군들이다.

2. 연산자들이 모든 자기동형이라 하자. 허용부분군은 모든 자기동형 아래에서 자기 자신으로 가는 \glqq{}특성부분군\grqq{}들이다.

3. \(\mG\)가 환, 곧 덧셈에 관해서는 아벨군이고 곱셈도 정의되어
\[
\begin{aligned}
r(a+b)&=ra+rb,\\
(a+b)r&=ar+br,\\
ab\cdot c&=a\cdot bc
\end{aligned}
\]
를 만족한다고 하자. 각 원소 \(r\)은 두 연산자 \(rx\), \(xr\)를 동시에 정의한다. 허용부분군은 \glqq{}아이디얼\grqq{} \(\mathfrak a\)이다. 연산 \(rx\)를 허용하여 \(r\mathfrak a\subseteq\mathfrak a\)인 것은 \emph{좌측아이디얼}, 연산 \(xr\)를 허용하여 \(\mathfrak a r\subseteq\mathfrak a\)인 것은 \emph{우측아이디얼}, 두 연산을 모두 허용하는 것은 \emph{양쪽아이디얼}이다. 환 \(\mG\)가 \emph{체}, 곧 \(\mG-\{0\}\)이 곱셈에 관하여 군이면 모든 아이디얼은 자명하다, 즉 \(\{0\}\) 또는 \(\mG\)이다\footnote{따라서 환에도 체에도 곱셈의 교환법칙을 가정하지 않는다.}.

4. \emph{환 \(\mo\)에 관한 가군.} \(\mo\)를 환, \(\mM\)을 덧셈으로 쓴 아벨군이라 하자. \(\mo\)의 원소와 \(\mM\)의 원소 사이에 곱 \(r\cdot a\)가 정의되어 다시 \(\mM\)에 들어 있고,
\[
\left.
\begin{aligned}
 r(a+b)&=ra+rb,\\
 rs\cdot a&=r\cdot sa,\\
 (r+s)a&=ra+sa
\end{aligned}
\right\}\qquad r,s\in\mo,\quad a,b\in\mM
\]
를 만족한다고 하자. 그러면 \(\mM\)을 \emph{\(\mo\)-가군}이라 한다. 더 정확히는 \(\mo\)의 곱셈자들을 왼쪽에 쓰므로 \emph{\(\mo\)-좌측가군}이다. 환 \(\mo\)는 동시에 연산자영역이다. 이 영역이 절대적, 곧 서로 다른 환원소들이 서로 다른 연산을 낳으면 \(\mo\)를 가군 \(\mM\)의 \emph{절대 곱셈자영역}이라 한다.

위와 같이 같은 연산을 생성하는 원소들을 같게 놓으면 임의의 곱셈자영역에서 절대 곱셈자영역으로 넘어갈 수 있다. 절대 곱셈자영역은 다시 환이다. 가군 \(\mM\)의 허용부분군을 \emph{부분가군}이라 한다. 특별히 \(\mo=\mM\)이면 좌측아이디얼로 돌아온다\footnote{같은 방식으로 우측가군에는
\[
\begin{aligned}
(a+b)r&=ar+br,\\
a\cdot rs&=ar\cdot s,\\
a(r+s)&=ar+as
\end{aligned}
\]
를 요구한다.}.

5. \emph{쌍가군.} \(\mM\)이 동시에 \(\mo\)-좌측가군과 \(\mo'\)-우측가군이고, 모든 \(a\in\mo\), \(\alpha\in\mM\), \(a'\in\mo'\)에 대해
\[
        a\cdot\alpha a'=a\alpha\cdot a'
\]
이면 \(\mM\)을 \emph{쌍가군}이라 한다.

6. \emph{아벨군의 자기동형환.}\footnote{A. Châtelet, \emph{Les Groupes Abéliens finis}, Paris, Gauthier-Villars, 1925, 99쪽을 보라.} 덧셈으로 쓴 아벨군에서 자기 자신으로 가는 준동형들에는
\[
        (H+\Theta)a=Ha+\Theta a,\qquad
        (H\Theta)a=H(\Theta a)
\]
로 덧셈과 곱셈을 정의할 수 있다. 이렇게 정의한 연산들도 분명 다시 준동형이고, 따라서 그 계에 속한다. 이 계는 환을 이루며, 4에 따라 그 아벨군을 이 환에 관한 가군으로 볼 수 있다.

이제부터 \glqq{}부분군\grqq{}은 언제나 허용부분군을 뜻한다. 두 허용부분군의 교집합 \(\mA\cap\mB\)은 다시 허용부분군이다. 두 부분군 가운데 적어도 하나가 정규부분군이면 그 곱 \(\mA\mB\)도 허용부분군이다.

\subsection*{§ 2. 동형정리.}

같은 연산자영역을 갖는 군 \(\mG,\overline{\mG}\) 사이의 사상이 첫째 통상적인 의미의 준동형 \(ab\mapsto\bar a\bar b\)이고, 둘째 \(a\)가 \(\bar a\)로 가면 \(\Theta a\)가 \(\Theta\bar a\)로 갈 때 이를 \emph{연산자준동형}이라 한다\footnote{연산자준동형의 개념은 \(\mG,\overline{\mG}\)가 서로 다른 연산자영역을 갖는 경우로 확장할 수 있다. 이때 준동형은 \(\mG\)를 \(\overline{\mG}\)로 사상할 뿐 아니라 연산자영역들도 서로 사상하며, \(a\)가 \(\bar a\)로, \(\Theta\)가 \(\overline\Theta\)로 가면 \(\overline\Theta\bar a\)가 \(\Theta a\)의 상이 되게 한다. 이 일반적 형태에서 연산자준동형은 환준동형도 포함한다. § 8을 보라.}. 기호는 \(\mG\sim\overline{\mG}\)이다. 대응이 일대일이면 \emph{연산자동형}이라 하고 \(\mG\simeq\overline{\mG}\)로 쓴다. 통상적인 군과 마찬가지로 다음 \emph{준동형정리}를 증명한다.

\emph{\(\overline{\mG}\)가 \(\mG\)의 연산자준동형상이면, \(\overline{\mG}\)는 한 몫군 \(\mG/\mN\)과 연산자동형이다. 여기서 \(\mN\)은 \(\overline{\mG}\)의 단위원에 대응하는 \(\mG\)의 모든 원소로 이루어진 허용 정규부분군이다. 역으로 모든 허용 정규부분군 \(\mN\)은 \(\mG\)의 연산자영역을 허용하고 \(\mG\)의 연산자준동형상인 몫군 \(\mG/\mN\)을 정의한다.}

자기 자신과 단위군 이외의 정규부분군을 갖지 않는 군을 \emph{단순}하다고 한다. 단순군의 모든 준동형은 동형이거나 모든 원소를 단위원으로 보낸다.

\emph{제1동형정리.} \emph{\(\overline{\mG}\)가 \(\mG\)의 준동형상이고 \(\overline{\mA}\)가 \(\overline{\mG}\)의 정규부분군이며, \(\mA\)가 \(\overline{\mA}\)의 원소들에 대응하는 \(\mG\)의 모든 원소의 전체이면 \(\mA\)도 정규부분군이고}
\[
        \overline{\mG}/\overline{\mA}\simeq\mG/\mA. \tag{1}
\]
\emph{증명.} \(\mG\sim\overline{\mG}\), \(\overline{\mG}\sim\overline{\mG}/\overline{\mA}\)이므로 \(\mG\sim\overline{\mG}/\overline{\mA}\)이다. 따라서 \(\overline{\mG}/\overline{\mA}\)는 \(\mG\)의 한 몫군과 동형이다. 대응 정규부분군은 \(\overline{\mG}/\overline{\mA}\)의 단위원에 대응하는 원소들의 전체이며, 이것이 바로 \(\mA\)이다.

\emph{덧붙임.} 준동형정리에 따라 \(\overline{\mG}=\mG/\mN\)으로 놓으면 \(\mA\)는 반드시 \(\mN\)을 포함한다. \(\mA\)에서 \(\overline{\mA}\)를 \(\overline{\mA}=\mA/\mN\)으로 되찾을 수 있고, 대응 \(\mA\leftrightarrow\overline{\mA}\)는 일대일이다. 식 (1)은
\[
        (\mG/\mN)/(\mA/\mN)\simeq\mG/\mA
\]
로도 쓸 수 있다.

\emph{제2동형정리.} \emph{\(\mA\)가 \(\mG\)의 부분군이고 \(\mB\)가 정규부분군이면 \(\mA\cap\mB\)는 \(\mA\)의 정규부분군이고}
\[
        \mA\mB/\mB\simeq\mA/(\mA\cap\mB).
\]
\emph{증명.} 준동형 \(\mG\sim\mG/\mB\)에서는 특히 \(\mA\)의 원소 \(a\)에 잉여류 \(a\mB\)가 대응하고, 이 잉여류들은 함께 군 \(\mA\mB/\mB\)를 이룬다. 따라서 \(\mA\sim\mA\mB/\mB\)이고, 준동형정리에서 주장이 따른다.

군 \(\mG\)에서 자기 자신 또는 한 부분군으로 가는 연산자준동형상을 \(\mG\)의 \emph{자기준동형}이라 한다. 특별히 \(\mG\)가 연산자를 갖는 덧셈 아벨군이고 § 1의 예 6에서처럼 준동형의 합과 곱을 정의하면, 군에서 자기 자신으로 가는 연산자준동형들은 \emph{자기동형환}이라는 한 환을 이룬다.

\emph{연산자를 갖는 단순 아벨군의 자기동형환은 체이다.}

\emph{증명.} 다른 허용부분군이 없으므로 모든 준동형은 군을 자기 자신으로 보내거나 영군으로 보낸다. 모든 것을 0으로 보내지 않는 준동형들은 위에서 단순군에 관해 말한 바에 따라 동형이고, 한 군에서 자기 자신으로 가는 동형사상들은 군을 이룬다. 따라서 환에서 영연산자를 빼면 군이 되므로 환 자체는 체이다.

특별히 \(\mG\)가 \(\mo\)-좌측가군이고 연산자준동형들을 우측연산자로 쓰면, \(\mG\)는 \emph{왼쪽의 \(\mo\)와 오른쪽의 자기동형환에 관한 쌍가군}이 된다. 실제로 새 연산자 \(\Gamma\)들을 준동형으로 택했다는 사실은
\[
        (a+b)\Gamma=a\Gamma+b\Gamma,
        \qquad
        ra\cdot\Gamma=r\cdot a\Gamma
\]
로 나타난다. 이 식들은 이미 앞서 해석한 나머지 식들과 함께 쌍가군을 특징짓는다.

역으로 한 쌍가군이 주어지면 바로 이 \emph{식들}은 우측 곱셈자영역의 각 원소가 좌측연산자들에 관한 자기준동형을 유도한다는 것과 그 역을 표현한다.

\subsection*{§ 3. 합성열.}

\(\mG\)에 허용부분군들의 유한열
\[
        \mG>\mA_1>\mA_2>\cdots>\mA_r=\mE \tag{2}
\]
이 존재하고 --- \(\mE\)는 단위원 \(e\)만으로 이루어진 군이다 --- 각 \(\mA_i\)가 앞선 군의 정규부분군이며, 열의 연속한 두 항 사이에 같은 성질을 가진 다른 항을 끼워 넣을 수 없으면 이 열을 \emph{합성열}이라 한다.

몫군
\[
        \mG/\mA_1,\quad \mA_1/\mA_2,\ldots,\quad \mA_{r-1}/\mE
\]
를 \emph{합성인자}라 한다. 이들은 \emph{단순군}, 곧 자기 자신과 단위군 이외의 허용 정규부분군을 갖지 않는 군이다.

연산자들 가운데 특히 모든 내부자기동형을 포함하면 열 (2)는 모두 정규부분군으로 이루어지고 이를 \emph{주열}이라 한다. 모든 자기동형을 포함하면 \emph{특성열}이라 한다. § 1의 뒤쪽 예들에서는 \(\mo\)의 아이디얼, 가군 또는 쌍가군의 합성열이 된다.

\emph{조르당--횔더 정리.} \emph{한 군 \(\mG\)에 서로 다른 두 합성열}
\[
\mG>\mA_1>\mA_2>\cdots>\mA_r=\mE,
\qquad
\mG>\mB_1>\mB_2>\cdots>\mB_s=\mE
\]
\emph{이 존재하면 길이는 같아 \(r=s\)이고, 합성인자들}
\[
        \mG/\mA_1,\quad \mA_1/\mA_2,\ldots,\quad\mA_{r-1}/\mE
\]
\emph{은 어떤 순서로}
\[
        \mG/\mB_1,\quad \mB_1/\mB_2,\ldots,\quad\mB_{s-1}/\mE
\]
\emph{와 동형이다.}

증명은 예컨대 E. Noether, \emph{Abstrakter Aufbau usw.}, Math. Annalen 96 (1926), § 10, 57쪽을 보라. 그곳에서는 동시에 다음도 증명한다.

\emph{\(\mG\)에 합성열이 하나 존재하면 \(\mG\)의 모든 정규부분군 \(\mH\)를 지나는 합성열을 그을 수 있다.}

합성열의 존재는 유한군에 대해 분명하며, 더 일반적으로 극대조건과 극소조건이 성립하는 군들에 대해서도 분명하다.

\emph{극대조건}은 부분군들의 모든 집합에 그 집합의 다른 부분군에 더 이상 포함되지 않는 극대부분군이 존재한다는 뜻이다. 이는 모든 상승사슬 \(\mA_1\subset\mA_2\subset\mA_3\cdots\)이 유한개의 항 뒤에 멈춘다는 것과 동치이다.

\emph{극소조건}은 부분군들의 모든 집합에 그 집합의 다른 부분군을 더 이상 포함하지 않는 극소부분군이 존재한다는 뜻이다. 또는 모든 하강사슬 \(\mA_1\supset\mA_2\supset\mA_3\cdots\)이 유한개의 항 뒤에 멈춘다는 뜻이다.

\emph{부분군으로 정규부분군만 허용하면}, 예컨대 아벨군의 경우에는 \emph{합성열의 존재에서 역으로 극대조건과 극소조건이 따른다.}

\emph{증명.} 모든 정규부분군은 합성열을 가지며, 그 길이를 그 정규부분군의 길이라고 한다. \(\mA\subset\mB\)이면 \(\mB\)의 합성열을 \(\mA\)를 지나도록 그을 수 있으므로 \(\mA\)의 길이는 \(\mB\)의 길이보다 작다. 따라서 최소길이의 모든 부분군은 극소이고 최대길이의 모든 부분군은 극대이다.

\subsection*{§ 4. 직접곱과 직접교집합.}

다음 조건들이 성립하면 군 \(\mG\)를 두 인자의 직접곱 \(\mG=\mA\times\mB\)라 한다.
\[
\begin{array}{ll}
1.& \mA,\mB\text{는 }\mG\text{의 정규부분군이다},\\
2.& \mA\mB=\mG,\\
3.& \mA\cap\mB=\mE.
\end{array}
\]
\emph{이 정의는 다음과 분명히 동치이다.} \(\mG\)의 모든 원소 \(g\)는 \(a\in\mA\), \(b\in\mB\)인 \(g=ab\)로 유일하게 표현되고, \(\mA\)의 원소들은 \(\mB\)의 원소들과 가환한다, 즉 \(ab=ba\)이다.

\(\mB_i=\mA_1\cdots\mA_{i-1}\mA_{i+1}\cdots\mA_n\)으로 놓을 때 모든 \(i\)에 대해 \(\mG=\mA_i\times\mB_i\)가 직접곱이면, 군 \(\mG\)를 \(n\)개 인자의 직접곱 \(\mG=\mA_1\times\cdots\times\mA_n\)이라 한다.

\emph{이 정의는 다시 다음과 동치이다.} \(\mG\)의 모든 \(g\)는
\[
        g=a_1a_2\cdots a_n,\qquad a_i\in\mA_i
\]
로 유일하게 표현되고, \(\mA_i\)의 원소들은 \(\mA_k\)의 원소들과 가환한다.

다음 사실들을 자주 쓴다.
\begin{enumerate}
\item \(\mA_1\times\cdots\times\mA_n=\mK\), \(\mK\times\mH=\mG\)이면 \(\mA_1\times\cdots\times\mA_n\times\mH=\mG\)이다.
\item \(\mG=\mA_1\times\cdots\times\mA_n=\mC_1\times\cdots\times\mC_n\)이고 \(\mC_i\subseteq\mA_i\)이면 \(\mC_i=\mA_i\)이다. 이는 \(\mA_i\)의 원소들을 \(\mC_j\)들로 나타내는 표현에서 따른다.
\item \(\mG=\mA\times\mB\), \(\mK\supseteq\mA\)이면 \(\mK=\mA\times(\mK\cap\mB)\)이다. 이는 \(\mK\)의 원소들을 \(ab\) 꼴로 나타낼 때 둘째 인자가 \(\mK,\mB\)에 모두 속한다는 데서 따른다.
\item \(\mG=\mA\times\mB\)이면 \(\mA\simeq\mG/\mB\), \(\mB\simeq\mG/\mA\)이다(제2동형정리). 여기서 다시 다음이 따른다.
\item \(\mG=\mA\times\mB\)이고 \(\mA,\mB\)가 각각 길이 \(m,n\)의 합성열을 가지면 \(\mG\)는 길이 \(m+n\)의 합성열을 가진다.
\item \(\mA\simeq\overline{\mA}\), \(\mB\simeq\overline{\mB}\)이면 \(\mA\times\mB\simeq\overline{\mA}\times\overline{\mB}\)이다.
\end{enumerate}

단위군 \(\mE\)와 다른 인자들의 직접곱으로 표현되지 않는 군을 직접분해불능이라 한다.

분명히 \emph{극소조건을 갖는 모든 군은 유한개의 직접분해불능군의 직접곱이다}\footnote{가환인 경우 W. Krull이, 일반적인 경우 O. Schmidt가 보였듯이(각주 6 참조), 극대조건과 극소조건 아래에서 이 표현은 동형을 제외하면 유일하다. 직접분해불능이라는 개념을 바꾸지 않고 모든 내부자기동형을 연산자영역에 더할 수 있으므로, 정규부분군에 대한 유한성 조건 또는 주열의 존재만을 요구해도 충분하다.}.

문제의 군들을 덧셈으로 쓰면 곱과 직접곱이라는 개념은 합과 직접합으로 바뀐다. 군들의 합은 \((\mA_1,\mA_2,\ldots)\), 직접합은 \(\mA_1+\mA_2+\cdots\)로 쓴다.

\glqq{}직접교집합\grqq{}이라는 개념은 뒤에서 쓰지 않지만, 직접교집합과 직접곱의 관계에 관한 다음 정리들은 다음 장의 아이디얼론을 합 대신 교집합으로도 정식화할 수 있음을 보여 준다. 이로써 알려진 가환이론들과의 관계가 세워진다.

다음 조건들이 성립하면 군 \(\mD\)를 \(\mG\)에 관한 두 군 \(\mR,\mS\)의 직접교집합이라 한다.
\[
\begin{array}{ll}
1.& \mR,\mS\text{는 }\mG\text{의 정규부분군이다},\\
2.& \mR\cap\mS=\mD,\\
3.& \mR\mS=\mG.
\end{array}
\]
마찬가지로 \(\mR_i=\mS_1\cap\cdots\cap\mS_{i-1}\cap\mS_{i+1}\cap\cdots\cap\mS_n\)으로 놓을 때 모든 \(i\)에 대해 \(\mD=\mR_i\cap\mS_i\)가 직접교집합이면, \(\mD\)를 \(n\)개 군 \(\mS_1,\ldots,\mS_n\)의 직접교집합이라 한다.

정의에 따라 \(\mD\)는 \(\mG\)의 정규부분군이어야 한다. \(\mD=\mS_1\cap\cdots\cap\mS_n\)이 직접교집합이고 준동형 \(\mG\sim\mG/\mD\)에서 군 \(\mG,\mS,\mD\)가 \(\overline{\mG},\overline{\mS},\overline{\mE}\)로 가면 \(\overline{\mE}=\overline{\mS}_1\cap\cdots\cap\overline{\mS}_n\)은 직접교집합이다. 역으로 그러한 모든 표현 \(\overline{\mE}=\overline{\mS}_1\cap\cdots\cap\overline{\mS}_n\)에서 직접표현 \(\mD=\mS_1\cap\cdots\cap\mS_n\)으로 돌아갈 수 있다. 이 대응으로 임의의 \(\mD\)에 대한 직접교집합 정리들은 특수한 경우 \(\mD=\mE\)로 환원된다.

\(\mD=\mE\)이고 인자가 둘이면 직접곱과 직접교집합의 조건은 일치한다. \(n\)개 인자에 대해서는 직접곱과 직접교집합 사이에 다음 일대일 관계가 있다.

\emph{I. \(\mG=\mA_1\times\cdots\times\mA_n=\mA_i\times\mB_i\)이면 \(\mE=\mB_1\cap\cdots\cap\mB_n=\mB_i\cap\mC_i\)는 직접교집합이고 \(\mC_i=\mA_i\)이다.}

\emph{II. \(\mE=\mS_1\cap\cdots\cap\mS_n=\mR_i\cap\mS_i\)가 직접교집합이면 \(\mG=\mR_1\times\cdots\times\mR_n=\mR_i\times\mT_i\)이고 \(\mT_i=\mS_i\)이다.}

\emph{I의 증명.} \(\mB=\mB_1\cap\cdots\cap\mB_n=\mB_i\cap\mC_i\)라 하자. 그러면 \(\mC_i\supseteq\mA_i\)이고, \(\mC_i\subseteq\mA_i\)임을 보이면 된다. 예컨대 \(c\in\mC_1\)이면 \(c\in\mB_2,\ldots,\mB_n\)이므로 \(\mA\)들에 의한 성분표현은
\[
        c=a_1a_2\cdots a_n
         =a_1ea_3\cdots a_n
         =a_1a_2e\cdots a_n
         =\cdots
         =a_1\cdots a_{n-1}e
\]
이다. 모든 \(\mA_i\)의 곱이 직접곱이므로 이 표현들은 일치한다. 첫 자리를 제외한 각 자리에는 한 번씩 \(e\)가 있으므로 \(c=a_1\), 곧 \(\mC_1\subseteq\mA_1\)이다. 따라서 \(\mC_1=\mA_1\)이고 마찬가지로 \(\mC_i=\mA_i\)이다. 이제 \(\mB=\mB_i\cap\mC_i=\mB_i\cap\mA_i=\mE\), \(\mB_i\mC_i=\mB_i\mA_i=\mG\)이므로 이는 직접교집합이다.

\emph{II의 증명.} \(\mH=\mR_1\cdots\mR_n=\mR_i\mT_i\)라 하자. 그러면 \(\mT_i\subseteq\mS_i\)이고, \(\mS_i\subseteq\mT_i\)임을 보이면 된다. 예컨대 \(s\in\mS_1\)이라 하자. \(\mG=\mS_i\times\mR_i\)에 따라
\[
        s=s_1e=s_2r_2=\cdots=s_nr_n.
\]
\(t_1=er_2\cdots r_n=t_ir_i\)로 놓으면, \(\mR_i,\mS_i\)의 원소별 가환성에 따라
\[
        st_1^{-1}=s_it_i^{-1}
\]
이다. 이는 모든 \(\mS_i\)의 원소이므로 \(e\)와 같다. 따라서 \(\mS_1\subseteq\mT_1\), 곧 \(\mT_1=\mS_1\)이고, 마찬가지로 \(\mT_i=\mS_i\)이다. 그러면 \(\mH=\mR_i\mT_i=\mR_i\mS_i=\mG\), \(\mR_i\cap\mT_i=\mR_i\cap\mS_i=\mE\)이므로 \(\mG\)는 \(\mR_i\)들의 직접곱이다.

\clearpage

\providecommand{\mC}{\mathfrak{C}}
\providecommand{\mD}{\mathfrak{D}}
\providecommand{\mK}{\mathfrak{K}}
\providecommand{\mR}{\mathfrak{R}}
\providecommand{\mS}{\mathfrak{S}}
\providecommand{\mT}{\mathfrak{T}}
\providecommand{\mO}{\mathfrak{O}}
\providecommand{\ma}{\mathfrak{a}}
\providecommand{\mb}{\mathfrak{b}}
\providecommand{\md}{\mathfrak{d}}
\providecommand{\mf}{\mathfrak{f}}
\providecommand{\mh}{\mathfrak{h}}
\providecommand{\mn}{\mathfrak{n}}
\providecommand{\mt}{\mathfrak{t}}
\providecommand{\mx}{\mathfrak{x}}
\providecommand{\mz}{\mathfrak{z}}
\providecommand{\mL}{\mathfrak{L}}
\providecommand{\iso}{\simeq}
\providecommand{\ann}{\operatorname{Ann}}
\providecommand{\tuple}[1]{(#1)}

\subsection*{§ 5. 완전가약군.}

한 군이 유한개의 단순군의 직접곱
\[
        \mG=\mA_1\times\cdots\times\mA_n
\]
이면 그 군을 \emph{완전가약}이라 한다. 이 경우 군들
\[
        \mA_1\times\cdots\times\mA_n
        >\mA_1\times\cdots\times\mA_{n-1}
        >\cdots>\mA_1>\mE
\]
은 합성열을 이룬다. 그 길이는 \(n\)이고 합성인자들은
\[
        \simeq\mA_n,\ \mA_{n-1},\ldots,\mA_1
\]
이다.

\subsection*{§ 6. 체에 관한 가군. 초복소계.}

한 체, 반드시 가환일 필요는 없는 체에 관한 유한기저 가군은 연산자를 갖는 완전가약군의 거의 자명한 예이다.

\(\mG\)를 \(K\)-우측가군이라 하고 \(K\)의 단위원 \(e\)가 동시에 단위연산자, 곧 모든 \(a\in\mG\)에 대해 \(ae=a\)라고 하자. 한 원소 \(a\)에서 생성되는 모든 부분가군 \(aK\)는 단순하다. 실제로 그 안의 임의의 비영원소에서 생성되는 가군과 같다. 따라서 \(\mH\)가 어떤 원소들에서 생성되는 부분가군이고 \(a\notin\mH\)이면 \(\mH\cap aK=\mE\), 따라서 \(\mH+aK\)는 직접합이다. 이처럼 한 기저원소 \(a_1\)에서 시작하여 새 기저원소 \(a_2,a_3,\ldots\)를 계속 더하면 \(\mG\)를 직접합
\[
        \mG=a_1K+a_2K+\cdots+a_nK
\]
으로 얻는다. \emph{따라서 \(\mG\)는 완전가약이다.}

합성열의 길이 \(n\)을 \(K\)에 관한 \(\mG\)의 \emph{랭크}라 한다. 기저 \(a_1,\ldots,a_n\)은 원소들의 표현이 유일하고 \(e\)를 단위연산자로 가정했으므로 선형독립이다.

모든 부분가군 \(\mH\)는 직접가수이다.
\[
        \mG=\mH+\mR
        =(h_1K+\cdots+h_sK)+(r_1K+\cdots+r_{n-s}K).
\]
달리 말해, \emph{\(\mH\)의 모든 선형독립기저를 \(\mG\)의 선형독립기저로 보충할 수 있다.} \(r_i\)들은 원래의 기저원소 \(a_i\)들 가운데서 택할 수도 있다.

\((a_1,\ldots,a_n)\), \((c_1,\ldots,c_n)\)을 선형독립기저라 하자. 그러면
\[
        c_k=\sum_i a_i\pi_{ik},
\]
또는 행렬로
\[
        (c_1,\ldots,c_n)=(a_1,\ldots,a_n)P,
        \qquad
        P=\begin{pmatrix}
        \pi_{11}&\cdots&\pi_{1n}\\
        \vdots&&\vdots\\
        \pi_{n1}&\cdots&\pi_{nn}
        \end{pmatrix}.
\]
역으로
\[
        (a_1,\ldots,a_n)=(c_1,\ldots,c_n)Q.
\]
따라서
\[
        (a_1,\ldots,a_n)=(a_1,\ldots,a_n)PQ,
        \qquad
        (c_1,\ldots,c_n)=(c_1,\ldots,c_n)QP.
\]
\(a_i,c_i\)들이 선형독립이므로
\[
        PQ=QP=E.
\]

동시에 환이기도 한 특수한 \(K\)-가군이 \glqq{}초복소계\grqq{}이다.

가환체 \(K\)에 관한 우측가군이기도 한 환 \(\mo\)가 다음 조건들을 만족하면 \emph{\(K\)에 관한 초복소계}, 문헌에서 말하는 \glqq{\(K\) 위의 대수\grqq{}}라 한다.
\begin{enumerate}
\item 랭크가 유한하다. 선형독립기저를 \(u_1,\ldots,u_n\)이라 하자.
\item \(ab\cdot x=a\cdot bx=ax\cdot b\)이다. 이를 \(K\)가 \(\mo\)와 가환적으로 결합되어 있다고 표현한다.
\item \(K\)의 단위원 \(e\)가 동시에 단위연산자이다. 즉 모든 \(a\in\mo\)에 대해 \(ae=a\)이다.\footnotemark
\end{enumerate}
\footnotetext{조건 2는 \(\mo\)를 \(\mo\)-좌측가군으로 보거나 \(\mo\)-우측가군으로 볼 때 \(K\)의 한 원소에 의한 곱셈이 \(\mo\)의 준동형이라는 뜻이다. 조건 3에 따라 이 준동형들은 실제로 동형이다.}

\[
        \Bigl(\sum_i u_i\alpha_i\Bigr)
        \Bigl(\sum_k u_k\beta_k\Bigr)
        =\sum_{i,k}(u_i u_k)(\alpha_i\beta_k)
\]
이므로, \(K\) 이외에 각 곱 \(u_i u_k\)를 \(u_j\)들로 나타내는 방법
\[
        u_i u_k=\sum_j u_j\gamma^j_{ik}
\]
을 주는 \emph{곱셈표}가 주어지면 계는 유일하게 결정된다. \(\gamma^j_{ik}\)들은 결합법칙에서 나오는 알려진 관계들을 만족해야 한다.

\(\mo\)가 단위원 \(e\), 곧 모든 \(a\)에 대해 \(ea=ae=a\)인 원소를 갖는다고 가정하는 경우가 많다. 그러면 \(K\)의 원소 \(x\)를 \(ex\)와 동일시하여 \(K\)를 \(\mo\)의 부분체로 볼 수 있다. 이때 초복소계는 \emph{중심에 놓인 한 체에 관하여 유한랭크인 환}이라고도 기술할 수 있다.

유한군의 원소들을 기저원소 \(u_i\)로 삼고 군의 곱셈을 곱셈으로 삼아 얻는 \emph{군환}은 초복소계의 한 예이다. 이때 체 \(K\)는 임의이다.

\subsection*{§ 7. 행렬.}

한 환 \(\mo\)의 원소 \(\pi_{ik}\)로 된 정방행렬들은 통상적인 행렬곱
\[
        \sum_j\pi_{ij}\rho_{jk}=\sigma_{ik}
\]
과 덧셈 \(\pi_{ik}+\rho_{ik}=\sigma_{ik}\)에 관하여 환을 이룬다.

체 \(K\)의 원소로 된 행렬이 우측역행렬과 좌측역행렬을 모두 가지면 \emph{정칙}이라 한다.

§ 6에서 보았다. \emph{\(K\)-가군의 두 선형독립기저 사이의 변환행렬은 정칙이다.}

\emph{정칙이기 위해서는 우측역행렬만으로 충분하다.} 실제로 선형독립기저 \((a_1,\ldots,a_n)\)을 갖는 \(K\)-우측가군, 곧 부정원 \(a_1,\ldots,a_n\)의 선형형식가군을 만들고
\[
        (c_1,\ldots,c_n)=(a_1,\ldots,a_n)P
\]
로 놓자. 그러면
\[
        (c_1,\ldots,c_n)P^{-1}
        =(a_1,\ldots,a_n)PP^{-1}
        =(a_1,\ldots,a_n).
\]
따라서 \((c_1,\ldots,c_n)\)은 다시 기저이고 선형독립이므로 행렬 \(P\)는 정칙이다. 또한 우측역행렬은 좌측역행렬과 같다.

마찬가지로 좌측역행렬만으로도 충분하다.

\emph{\(P\)가 좌측역행렬을 가지면 \(P\)는 좌측 영인자가 아니다.}
\emph{증명.} \(PA=0\)에서 \(A=P^{-1}PA=0\)이 따른다.

따라서 \(P\)는 우측역행렬도 하나만 가지며, § 6에 따라 그것은 좌측역행렬과 일치한다.

통상대로 \(P_x\)가 행렬 \((\pi_{ik}x)\)를 뜻하고 \(C_{ik}\)가 \(ik\) 자리에는 단위원, 나머지 자리에는 0을 갖는 행렬이라 하면
\[
        (\pi_{ik})=\sum_{i,k}C_{ik}\pi_{ik}.
\]
따라서 \(K\)의 행렬들은 랭크 \(n^2\)인 \(K\)-가군을 이룬다. \(C_{ik}\)들은
\[
        C_{ij}C_{jk}=C_{ik},
        \qquad
        C_{ij}C_{\bar j k}=0,\quad j\ne\bar j
\]
를 만족한다. 그러므로 \(K\)의 행렬들은 유한 \(K\)-가군이고, \(K\)가 가환이면 초복소계를 이룬다.

\section*{제II장. 비가환 아이디얼론.}

\subsection*{§ 8. 환의 준동형정리.}

\(\ma\)가 \(\mo\)의 양쪽아이디얼이면 잉여류영역 \(\mo/\ma\)는 \(\mo\)-가군일 뿐 아니라 환이기도 하다는 것은 쉽게 알 수 있다. \(\mo\)의 모든 준동형상은 다시 환이며, \(\ma\)가 양쪽아이디얼인 잉여류환 \(\mo/\ma\)와 동형이다. 증명은 군의 경우와 같다.

특히 \(\mo\)가 체이면 \((0)\)과 \(\mo\) 이외의 아이디얼은 없다. 따라서 이 경우 준동형은 동형이거나 모든 원소를 0에 대응시킨다.

특히 환 \(\mo\)가 \(K\)-우측가군이고, \(K\)가 원소별로 \(\mo\)와 가환하는 환이라고 하자.
\[
        ab\cdot x=a\cdot bx=ax\cdot b
\]
(예컨대 가환체 \(K\)에 관한 초복소계를 다루는 경우가 그렇다.) 그러면 허용되는 우측ㆍ좌측ㆍ양쪽아이디얼은 동시에 \(K\)-가군인 것들뿐이며, 환준동형에 더하여 \(K\)와의 곱셈에 관한 연산자준동형도 요구한다. 허용아이디얼은 \(\mo\)와 \(K\)에 관한 쌍가군이 된다. (초복소계에서 단순히 \glqq{}아이디얼\grqq{}이라 할 때에는 언제나 바탕 계수체 \(K\)에 관하여 허용되는 것만을 뜻한다.) 이 확장에서도 위의 준동형정리는 성립한다.

환준동형의 한 예는 승수영역 \(\mo\)에서 절대승수영역으로 넘어가는 것이다(§ 1, 예 4). 따라서 절대승수영역은 잉여류환 \(\mo/\md\)와 동형이며, 여기서 \(\md\)는 \(\mM\)을 소멸시키는 \(\mo\)의 원소들이 이루는 양쪽아이디얼이다.

준동형정리가 성립하므로 § 2에서와 같이 환동형과 연산자동형에 관한 제1ㆍ제2 동형정리가 따른다.

\(\mo\)의 부분집합들의 모든 곱 \(\mA\mB\)는 가군곱으로 이해한다. 곧 \(\mA\mB\)는 \(a\in\mA\), \(b\in\mB\)인 모든 합 \(\sum ab\)의 전체이고, \(a\mB\)는 \(b\in\mB\)인 모든 합 \(\sum ab\)의 전체이다. \(\mB\)가 어떤 환을 연산자영역으로 하는 우측가군이면 곱 \(\mA\mB\), 또는 \(a\mB\)도 우측가군이다. (특히 우측아이디얼, 계수영역 \(K\)에 관한 허용아이디얼 등이다.) \((\mA,\mB)\)가 가군 \(\mA,\mB\)의 합을 나타낼 때 다음 계산법칙이 성립한다.
\[
        \mA\mB\cdot\mC=\mA\cdot\mB\mC,
\]
\[
        \mA(\mB,\mC)=(\mA\mB,\mA\mC),
        \qquad
        (\mA,\mB)\mC=(\mA\mC,\mB\mC).
\]
직접합은 \(\mA+\mB\)로 나타낸다(§ 4).

이하에서는 우측아이디얼에 관하여 극대조건 또는 극소조건(§ 3)을 만족하는 환을 자주 고찰한다. 특히 초복소계가 여기에 속한다. 위의 약속에 따라 \(K\)-가군만 아이디얼로 허용되며, 따라서 그 랭크가 유계이기 때문이다(§ 6).

\subsection*{§ 9. 멱등원. 우측아이디얼로의 직접합분해.}

\(\mo\)를 단위원을 갖는 환이라 하자.

\(\mr\)가 우측아이디얼이면 \(\mr\mo\subseteq\mr\)이고, 단위원 때문에 사실 \(\mr\mo=\mr\)이다. 좌측아이디얼에 대해서도 마찬가지로 \(\mo\ml=\ml\)이다.

원소 \(c\)가 \(c^2=c\)를 만족하면 \emph{멱등원}이라 한다.

\emph{\(\mo=\mr_1+\cdots+\mr_n\)이 우측아이디얼들로의 분해이고}
\[
        e=e_1+\cdots+e_n,
        \qquad e_i\in\mr_i
\]
\emph{이면}
\[
        e_i^2=e_i,
        \qquad e_i e_k=0\quad(i\ne k)
        \quad\hbox{(\glqq{}직교관계\grqq{})},
\]
\[
        \mr_i=e_i\mo
\]
\emph{이다.}

\emph{증명.} \(r\)을 \(\mr_1\)의 한 원소라 하자. 그러면
\[
        r=er=e_1r+\cdots+e_nr,
        \qquad e_ir\in\mr_i
\]
이지만, 한편
\[
        r=r+0+\cdots+0
\]
이기도 하므로
\[
        e_1r=r,
        \qquad e_ir=0\quad(i\ne1)
\]
이다. 첫째 식에서 \(\mr_1\subseteq e_1\mo\)가 따르고, 다른 한편 \(e_1\mo\subseteq\mr_1\)이므로 \(e_1\mo=\mr_1\)이다. \(r=e_1\)로 특수화하면
\[
        e_1^2=e_1,
        \qquad e_i e_1=0\quad(i\ne1)
\]
을 얻는다. 지수 1을 \(k\)로 바꾸어도 같은 결과가 성립한다.

\emph{역으로 \(e=\sum e_i\), \(e_i e_k=0\) \((i\ne k)\), \(e_i^2=e_i\)이고 \(\mr_i=e_i\mo\)로 놓으면 \(\mo=\mr_1+\cdots+\mr_n\)이다.}

\emph{증명.} \(\mo\)의 모든 원소 \(r\)은
\[
        r=er=e_1r+\cdots+e_nr
\]
이므로
\[
        \mo=(e_1\mo,\ldots,e_n\mo)
\]
이다. 그런데 이 표시는 유일하다. 실제로
\[
        0=e_1a_1+\cdots+e_na_n
\]
에서 \(e_i\)를 곱하면
\[
        e_i^2a_i=e_i a_i=0
\]
이 따르기 때문이다.\footnote{또 다른 역도 성립한다. \(e_i e_k=0\), \(e_i^2=e_i\), \(c=\sum e_i\)이면 \(c\)는 환 \(e_1\mo+\cdots+e_n\mo\)의 좌측단위원이다. 합이 직접합이라는 것은 위와 아주 똑같이 알 수 있다.}

따라서 우측표현
\[
        \mo=\mr_1+\cdots+\mr_n=e_1\mo+\cdots+e_n\mo
\]
에서 좌측표현
\[
        \mo=\ml_1+\cdots+\ml_n=\mo e_1+\cdots+\mo e_n
\]
이 따른다.

\(\mr_i\)들이 직접분해불능이면 \(\ml_i\)들도 그러하다. 실제로 예컨대 \(\ml_1\)의 분해가 있다면 이는
\[
        \ml_1=\mo e_1=\mo e_1'+\mo e_1'',
        \qquad
        \mo=\mo e_1'+\mo e_1''+\mo e_2+\cdots+\mo e_n
\]
을 뜻한다. 여기서 우측분해
\[
        \mo=e_1'\mo+e_1''\mo+e_2\mo+\cdots+e_n\mo
        =\mr_1'+\mr_1''+\mr_2+\cdots+\mr_n,
\]
\[
        \mr_1\simeq \mo/(\mr_2+\cdots+\mr_n)=\mr_1'+\mr_2'
\]
을 얻으므로 \(\mr_1\)은 분해가능하게 된다.

\subsection*{§ 10. 양쪽아이디얼로의 분해.}

\(\mo\)가 다시 단위원을 갖는 환이고
\begin{equation}
        \mo=\ma_1+\cdots+\ma_n                                  \tag{1}
\end{equation}
이 \(\mo\)를 양쪽아이디얼들로 분해한 것이면, 아이디얼들에 대해서도 직교관계
\begin{equation}
        \ma_i\ma_k=0\quad(i\ne k),
        \qquad
        \ma_i^2=\ma_i                                             \tag{2}
\end{equation}
가 성립한다.

\emph{증명.} \(\mb_1=\ma_2+\cdots+\ma_n\), \(\mo=\ma_1+\mb_1\)이라 하자. 그러면 \(\ma_1\mb_1\leq \ma_1\cap\mb_1=0\), 따라서 \(\ma_1\mb_1=0\)이고, 더구나
\[
        \ma_1\ma_i=0\quad(i\ne1)
\]
이다. 또한
\[
        \ma_1=\ma_1\mo=\ma_1(\ma_1+\mb_1)
        =(\ma_1^2,\ma_1\mb_1)=\ma_1^2.
\]

역으로, \emph{(1)과 (2)에서 가군 \(\ma_i\)들이 양쪽아이디얼이라는 것이 따른다.}
\emph{증명.}
\[
        \mo\ma_i=(\ma_1+\cdots+\ma_n)\ma_i=\ma_i^2=\ma_i,
        \qquad
        \ma_i\mo=\ma_i.
\]

\emph{환 \(\ma_i\)의 우측아이디얼은 \(\mo\)의 우측아이디얼이다.}
\emph{증명.}
\[
        \mr\mo=\mr(\ma_i+\mb_i)=(\mr\ma_i,\mr\mb_i)
        \subseteq(\mr,\ma_i\mb_i)=(\mr,0)=\mr.
\]

\emph{\(\mr\)가 \(\mo\)의 우측아이디얼이면 \(\mr\)는 \(\mr\ma_i\leq\ma_i\)인 우측아이디얼들의 직접합이다.}
\emph{증명.} \(\mr=\mr\mo=(\mr\ma_1,\ldots,\mr\ma_n)\)이고 \(\mr\ma_i\mo=\mr\ma_i\)이므로 \(\mr\ma_i\)는 \(\mr\)에 포함되는 우측아이디얼이다. 또한 \(\mr\ma_i\leq\ma_i\)이므로 합은 직접합이다.
\[
        \mr=\mr\ma_1+\cdots+\mr\ma_n.
\]

\emph{특히 \(\mr\)가 직접분해불능이면 \(\mr\)는 성분을 하나만 가질 수 있으므로 어떤 \(\ma_i\) 안에 들어가야 한다.}

이 정리들에 따라 각각의 \(\ma_i\)의 아이디얼론을 알면 \(\mo\)의 아이디얼론을 다룰 수 있다. 그러므로 대개 양쪽으로 분해불능인 환에 국한한다.

\emph{서로 다른 \(\ma_i\)에 들어 있는 두 우측아이디얼은 결코 연산자동형일 수 없다.}
\emph{증명.} \(\ma_1\) 안에 있는 아이디얼은 다른 모든 \(\ma_k\)에 의해서는 소멸되지만 \(\ma_1\)에 의해서는 소멸되지 않는다. 반면 \(\ma_2\) 안에 있는 아이디얼은 \(\ma_1\)에 의해 소멸된다.

따라서 환 \(\mo\)를 양쪽으로 분해불능인 아이디얼들의 합 \(\ma_1+\cdots+\ma_n\)으로 분해할 수 있고, 각 \(\ma_i\)를 다시 한쪽으로 분해불능인 우측아이디얼들의 합 \(\mr_i'+\mr_i''+\cdots\)로 분해할 수 있다면, 서로 연산자동형인 \(\mr\)들은 같은 \(\ma_i\)에 속하는 것들 사이에서 찾아야 한다. 그러나 다음 예가 보이듯 같은 \(\ma_i\) 안에도 서로 다른 \(\mr\)의 류, 곧 연산자동형이 아닌 류들이 있을 수 있다.

\(K\)를 유리수체라 하고
\[
        \mo=e_1K+e_2K+uK
\]
를 초복소계라 하자. \(e_i\)와 \(u\)는 \(K\)의 수들과 가환하며, 곱셈표는
\[
\begin{array}{c|ccc}
     & e_1&e_2&u\\ \hline
 e_1 & e_1&0&u\\
 e_2 & 0&e_2&0\\
 u   & 0&u&0
\end{array}
\]
이다. 단위원은 \(e=e_1+e_2\)이다. \(e_i\)들은 직교관계를 만족한다. 따라서 우측분해는
\[
        \mo=e_1\mo+e_2\mo=(e_1,u)+(e_2)
\]
이고, 좌측분해는
\[
        \mo=\mo e_1+\mo e_2=(e_1)+(e_2,u)
\]
이다. \(e_2\mo\)와 \(\mo e_1\)은 분해불능임이 명백하므로 \(\mo e_2\)와 \(e_1\mo\)도 분해불능이어야 한다.

양쪽분해는 불가능하다. 그러한 분해가 있다면 한 성분은 \(e_1\mo\)를, 다른 성분은 \(e_2\mo\)를 포함해야 한다. 그런데 첫째 성분은 \(u\)를 포함하고, 둘째 성분도 \(ue_2=u\)이므로 \(u\)를 포함한다.

그러나 \(e_1\mo\)와 \(e_2\mo\)는 \(K\)에 관한 랭크가 서로 다르므로 연산자동형이 아니다.

\emph{\(\mo=\ma_1+\cdots+\ma_n\)이고 \(\ma_i\)들이 양쪽으로 분해불능이면, 이들은 유일하게 결정된다.}

\emph{증명.} 또한 \(\mo=\mc_1+\cdots+\mc_n\)이라 하자. 그러면
\[
        \ma_i=\ma_i\mo=(\ma_i\mc_1,\ldots,\ma_i\mc_n).
\]
마지막 합은 \(\ma_i\mc_k\leq\mc_k\)이고 동시에 \(\ma_i\mc_k\leq\ma_i\)이므로 직접합이다. 그런데 \(\ma_i\)들이 분해불능이므로 \(\ma_i\mc_k\)는 하나, 곧 \(\ma_i\mc_{j_i}\)를 제외하면 모두 0이어야 한다. 따라서
\[
        \ma_i=\ma_i\mc_{j_i}.
\]
마찬가지로
\[
        \mc_{j_i}=\mo\mc_{j_i}
        =\ma_1\mc_{j_i}+\cdots+\ma_n\mc_{j_i},
\]
이고 \(\ma_i\mc_{j_i}\ne0\)이므로 나머지는 모두 0이며,
\[
        \mc_{j_i}=\ma_i\mc_{j_i}=\ma_i
\]
이다. 증명 끝.

\subsection*{§ 11. 중심.}

환 \(\mo\)의 중심 \(\mZ\)는 모든 환원소와 가환하는 \(z\)들의 전체이다(모든 \(a\)에 대해 \(za=az\)). \(\mZ\)는 가환환이다.

\(\mo\)의 모든 한쪽아이디얼 \(\ma\)는 \(\mZ\) 안에 \glqq{}수축아이디얼\grqq{} \(\ma\cap\mZ\)을 갖는다. 실제로 \(\ma\)와 \(\mZ\)가 모두 \(\mZ\)-가군이므로 \(\ma\cap\mZ\)도 그러하다.

\(\mZ\)의 모든 아이디얼 \(\mA\)는 확장아이디얼, 곧 \(\mA\)가 \(\mo\)에서 생성하는 아이디얼 \(\ma=(\mA,\mo\mA)\)를 갖는다. 이것은
\[
        \ma\mo=(\mo\mA,\mA)\mo=(\mo\mA\mo,\mA\mo)
        =(\mo^2\mA,\mo\mA)\subseteq\mo\mA\subseteq\ma
\]
이므로 양쪽아이디얼이다. \(\mo\)에 단위원이 있다고 가정하면 \(\ma=\mo\mA\)로 쓸 수 있다. 이 경우 다음 정리가 성립한다.

\emph{\(\mo\)의 모든 양쪽분해는 중심의 한 분해와 일대일로 대응한다. 즉}
\begin{equation}
        \mo=\ma_1+\cdots+\ma_n,
        \qquad \mA_i=\ma_i\cap\mZ                         \tag{1}
\end{equation}
\emph{에서}
\begin{equation}
        \mZ=\mA_1+\cdots+\mA_n,
        \qquad \mo\mA_i=\ma_i                              \tag{2}
\end{equation}
\emph{가 따르며, 역도 성립한다.}

\emph{증명.} \(z\)를 중심원소라 하고
\begin{equation}
        z=z_1+\cdots+z_n                                    \tag{3}
\end{equation}
을 (1)에 따른 그 분해라 하자. 그러면 \(z_i\)들은 중심원소이다. 실제로
\[
        za=z_1a+\cdots+z_na=az=az_1+\cdots+az_n
\]
이고, 합이 직접합이며 \(\ma_i\)들이 양쪽아이디얼이므로
\[
        z_i a=az_i
\]
이다. 따라서 \(z_i\in\ma_i\cap\mZ=\mA_i\)이고, (3)은 주장한 합분해를 뜻한다. 특히 \(e=\sum e_i\)가 되므로 \(e_i\)들은 중심원소이다. 끝으로
\[
        z=ze_1+\cdots+ze_n
\]
이므로
\[
        \mA_i=\mZ e_i,
        \qquad
        \mo\mA_i=\mo\mZ e_i=\mo e_i=\ma_i.
\]

역으로 \(\mZ=\mA_1+\cdots+\mA_n\)이 중심의 한 분해이면 § 9에 따라
\[
        \mA_i=\mZ e_i,
        \qquad e_i^2=e_i,
        \qquad e_i e_k=0\quad(i\ne k).
\]
따라서 직교관계와 § 9에 의해
\[
\begin{aligned}
        \mo&=\mo e=\mo e_1+\cdots+\mo e_n
        =\mo\mZ e_1+\cdots+\mo\mZ e_n\\
        &=\mo\mA_1+\cdots+\mo\mA_n
        =\ma_1+\cdots+\ma_n
\end{aligned}
\]
이다. 이제 \(\ma_i\cap\mZ=\mA_i'\)라 놓으면 주장의 첫째 부분에서
\[
        \mZ=\mA_1'+\cdots+\mA_n'\quad\hbox{(직접합)}
\]
임을 안다. 그런데 \(\mA_i'\supseteq\mA_i\)이므로 \(\mA_i'=\mA_i\)이다(§ 4, 2).

\emph{따름정리.} \(\ma_i\)와 \(\mA_i\)는 동시에 양쪽으로 직접분해가능하거나 동시에 그렇지 않다.

\subsection*{§ 12. 멱영아이디얼.}

아이디얼 \(\mc\)가 \(\mc^\rho=0\)을 만족하면 \emph{멱영}이라 한다.

\emph{예.} § 10의 아이디얼 \((u)\).

\emph{\(\mo\)에 멱영 우측아이디얼 \(\mr\)가 있으면 멱영 좌측아이디얼, 나아가 양쪽아이디얼도 있다.}

\emph{증명.} \(\mr^\rho=0\)이라 하자. 그러면 \(\mo\mr\)는 좌측아이디얼이다. \(\mo\)가 단위원을 포함하지 않아 \(\mo\mr\)가 0일 수도 있는 경우에는 \((\mo\mr,\mr)\)가 좌측아이디얼이며,
\[
        (\mo\mr)^\rho=\mo\mr\mo\mr\cdots\mo\mr
        \subseteq\mo\mr\mr\cdots\mr=\mo\mr^\rho=0.
\]

\emph{두 멱영 우측아이디얼의 합은 다시 멱영 우측아이디얼이다.}

\emph{증명.} \(\mc^\rho=0\), \(\md^\sigma=0\)이라 하자.
\[
        (\mc,\md)^{\rho+\sigma-1}
        =(\ldots,\md\mc\cdots\md\cdots\mc,\ldots)
\]
의 각 항에는 \(\mc\) 인자가 적어도 \(\rho\)개 있거나 \(\md\) 인자가 적어도 \(\sigma\)개 있다. 첫째 경우에는
\[
        \md\mc\cdots\md\cdots\mc\cdots
        \subseteq\md\mc\mc\cdots\mc=\md\mc^\rho=0
\]
이고, 둘째 경우에는 \(\md\) 인자들에 관하여 마찬가지이다. 따라서
\[
        (\mc,\md)^{\rho+\sigma-1}=0.
\]

이제 \(\mo\)의 우측아이디얼들에 대한 극대조건이 성립하면 극대한 멱영 우측아이디얼이 있다. 이것은 다른 모든 멱영 우측아이디얼을 포함한다. 그렇지 않으면 그러한 아이디얼 하나와 합을 만들어 더 큰 것을 얻을 수 있기 때문이다. 이를 \(\mc\)라 하자. \(\mo\mc\)도 멱영 우측아이디얼이므로 \(\mo\mc\subseteq\mc\)이고, 따라서 \(\mc\)는 좌측아이디얼이다. 모든 멱영 좌측아이디얼 \(\md\)도 \(\mc\)에 포함된다. 실제로 \((\md,\md\mo)\)는 멱영 우측아이디얼이므로 \(\md\subseteq\mc\)이다. \emph{따라서 다른 모든 우측 및 좌측 멱영아이디얼을 포함하는 극대 멱영 양쪽아이디얼 \(\mc\), 곧 \glqq{}근기\grqq{}가 있다.}

§ 10의 예에서는 \((u)\)가 극대 멱영아이디얼이다.

근기가 영아이디얼이면 \glqq{}근기 없는 환\grqq{}(\glqq{}반단순환\grqq{}, \glqq{}데데킨트계\grqq{})이라 한다. 근기에 따른 잉여류환은 언제나 근기 없는 환이다.

\emph{\(\mZ\)가 \(\mo\)의 중심이고 \(\mc\)가 \(\mo\)의 근기이면, \(\mC=\mc\cap\mZ\)는 \(\mZ\)의 근기이다.}

\emph{증명.} \(\mC\)가 멱영이라는 것은 명백하다. \(\mZ\)에 이를 진정하게 포함하는 멱영아이디얼 \(\overline{\mC}\)가 있다면, 이것을 \(\mc\)에 완전히 포함되지 않는 멱영아이디얼 \(\overline{\mc}\)로 확장할 수 있다.
\[
        (\overline{\mC}\mo)^\rho
        =\overline{\mC}^{\rho}\mo^\rho=0.
\]

\emph{따름정리.} \(\mo\)가 근기 없는 환이면 \(\mZ\)도 그러하다. 그러나 § 10의 예가 보이듯 그 역은 성립하지 않는다. \(\mo/\mc\)의 중심은 \(\mZ/\mC\)와 동형인 환을 포함하지만, 이것은 진부분환일 수 있다(§ 10의 예).

\subsection*{§ 13. 완전가약환.}

§ 5에 따라, 한 환이 유한개의 단순 우측아이디얼의 직접합이면 \emph{우측 완전가약}이라 한다.

\emph{단위원을 갖는 우측 완전가약환은 \((0)\)이 아닌 멱영아이디얼을 갖지 않으며, 따라서 근기를 갖지 않는다.}

\emph{증명.} 모든 우측아이디얼 \(\mr\)는 직접가수이다(§ 5).
\[
        \mo=\mt+\mr=e_1\mo+e_2\mo.
\]
\(e_2^2=e_2\)이므로 \(e_2^\rho=e_2\)이기도 하다. \(\mr\)가 멱영이라면 \(e_2^\rho=0\), 따라서 \(e_2=0\), 따라서 \(\mr=0\)이다. 증명 끝.

이제 두 가지 역을 보이자. 첫째, \emph{근기가 없고 우측아이디얼들에 대한 극소조건을 만족하는 환은 우측아이디얼에 관하여 완전가약이다.}

\emph{증명.} \(\mt\ne0\)을 극소 우측아이디얼이라 하자. \(\mt\)가 멱등원을 포함함을 보이겠다.

\(\mt^2\subseteq\mt\)이고 \(\mt^2\ne0\)이므로 \(\mt^2=\mt\)이다. \(\mt^2\)도 우측아이디얼이기 때문이다. 따라서 \(a\mt\ne0\)인 \(a\in\mt\)가 있다. 그런데 \(a\mt\)도 우측아이디얼이므로 \(a\mt=\mt\)이어야 한다.

\(a\)에 의해 소멸되는 \(\mt\)의 모든 \(b\), 곧 \(ab=0\)인 \(b\)들의 전체는 우측아이디얼이고 \(\mt\)와 같지 않으므로 0이다. 따라서 \(ab=0\)이면 \(b=0\)이다.

\(\mt=a\mt\)이므로 \(a\)는 \(ac\)의 꼴로 나타난다. 곧 \(a=ac\), \(c\ne0\)이다. 따라서
\[
        ac=ac^2,
        \qquad a(c-c^2)=0,
        \qquad c-c^2=0,
        \qquad c^2=c.
\]
이제 \(\mt\)가 직접가수임을 보이자. \(c\mo\)는 \(\mt\)에 포함되는 우측아이디얼이고, \(c^2=c\in c\mo\)이므로 0이 아니다. 따라서 \(c\mo=\mt\)이다. \(\mo\)의 모든 원소 \(r\)은
\[
        r=cr+(r-cr)
        \qquad\hbox{(한쪽 페어스 분해)}
\]
로 나타난다. \(cr\)인 원소들은 아이디얼 \(\mt\)를 이루고, \(r-cr\)인 원소들은 \(c\)에 의해 소멸되는 다른 우측아이디얼 \(\mr\)를 이룬다. 따라서
\[
        \mo=(\mt,\mr).
\]
\(\mt\)의 원소들은 \(c\)에 의해 소멸되지 않으므로 합은 직접합이다.
\begin{equation}
        \mo=\mt+\mr.                                             \tag{1}
\end{equation}

이제 극소조건에 따라 \(\mo\)를 직접분해불능인 아이디얼들의 직접합
\[
        \mo=\mr_1+\cdots+\mr_n
\]
으로 나타내면 \(\mr_i\)들은 단순, 곧 극소이어야 한다. 실제로 예컨대 \(\mr_1\)이 단순하지 않고 \(\mt\)가 \(\mr_1\) 안의 극소아이디얼이면 (1)에서
\[
        \mr_1=\mt+\mr\cap\mr_1\qquad\hbox{(§ 4, 3)}
\]
을 얻는다. 그러면 \(\mr_1\)이 분해가능하게 되어 모순이다. 따라서 \(\mo\)는 완전가약이다.

둘째, \emph{근기가 없고 극소조건을 만족하는 환은 단위원을 갖는다.}

먼저 좌측단위원의 존재를 보이기 위해서는 다음을 보이면 충분하다. \emph{우측아이디얼 \(\ma=c_1\mo\ne\mo\)가 좌측단위원 \(c_1\)을 가지면, 더 큰 길이를 가지면서 역시 좌측단위원 \(c\)을 갖는 우측아이디얼 \(c\mo\)가 존재한다.} 실제로 완전가약성, 따라서 모든 길이의 유계성은 이미 보였으므로 이 과정을 유한번 거치면 \(\mo\) 자체의 좌측단위원에 이른다.

그러므로 \(\ma=c_1\mo\), \(c_1^2=c_1\)이라 하자. 식
\[
        r=c_1r+(r-c_1r)
\]
에 의해 위에서와 같이 페어스 분해 \(\mo=c_1\mo+\mr\)가 주어진다. 위에서와 같이 \(c_2\)를 \(\mr\) 안의 멱등원이라 하자. 그러면 \(c_2^2=c_2\), \(c_1c_2=0\)이다. \(c_1\)이 \(\mr\)의 모든 원소를 소멸시키기 때문이다. 이제 \(e_1=c_1\), \(e_2=c_2-c_2c_1\)로 놓으면
\[
        e_1^2=e_1,
        \quad e_1e_2=0,
        \quad e_2e_1=0,
        \quad e_2c_2=c_2,
        \quad e_2\ne0,
        \quad e_2^2=e_2.
\]
따라서 각주 13의 논평에 따라 \(c=e_1+e_2\)는 환
\[
        c\mo=e_1\mo+e_2\mo
\]
의 좌측단위원이다. 이 환은 \(e_2^2=e_2\ne0\)이므로 우측아이디얼로서 \(\ma\)보다 더 긴 길이를 갖는다.

구성한 좌측단위원 \(e\)가 우측단위원이기도 함을 보이기 위해 좌측아이디얼들에 대한 페어스 분해
\[
        \mo=\mo e+\ml
\]
를 만들자. \(\ml e=0\), \(\ml=e\ml\)이므로 \(\ml^2=\ml e\ml=0\)이다. 가정에 따라 멱영 좌측아이디얼은 존재할 수 없으므로 \(\ml=0\)이다. 따라서 \(e\)는 우측단위원이기도 하며, 결국 단위원이다.

\emph{정리. 우측 완전가약성과 단위원의 존재에서 양쪽 완전가약성이 따른다.}
\[
        \mo=\md_1+\cdots+\md_s,
        \qquad \md_i\ \hbox{는 양쪽단순이다.}
\]

앞의 증명에서와 같이, 모든 극소 양쪽아이디얼 \(\ma\)가 직접가수임을 보이면 충분하다. \(\ma\)는 우측아이디얼로서 직접가수이다.
\[
        \mo=\ma+\mr=e_1\mo+e_2\mo.
\]
이에 대응하는 좌측분해는
\[
        \mo=\mc+\ml=\mo e_1+\mo e_2
\]
이다. 따라서
\[
        \mr\mc\subseteq\mr\ma\leq\mr\cap\ma=0,
        \qquad
        \ml\ma=\mo e_2e_1\mo=0,
\]
\[
        (\ma\ml)^2=\ma\ml\ma\ml=0,
        \qquad\hbox{따라서 }\ma\ml=0,
\]
\[
        \mr=\mr\mo=(\mr\mc,\mr\ml)=\mr\ml,
        \qquad
        \ml=\mo\ml=(\ma\ml,\mr\ml)=\mr\ml.
\]
그러므로 \(\mr=\ml\)이고, 따라서 \(\mr\)는 양쪽아이디얼이다. 이에 따라 \(\ma\)는 양쪽 직접가수이다. 그 뒤는 한쪽 완전가약성의 증명과 같다.

\(\md_i\)들은 환으로서도 양쪽단순이다. \(\md_i\)의 모든 아이디얼은 \(\mo\)의 아이디얼이기 때문이다. 이제 그 구조를 조사한다.

\subsection*{§ 14. 단위원을 갖는 양쪽단순 완전가약환.}

\(\mo\)를 그러한 환이라 하자.
\[
        \mo=\mr_1+\cdots+\mr_n=e_1\mo+\cdots+e_n\mo,
        \qquad \mr_i\ \hbox{는 단순},
\]
이면
\[
        \mo=\ml_1+\cdots+\ml_n=\mo e_1+\cdots+\mo e_n,
        \qquad \ml_i\ \hbox{는 직접분해불능}\quad\hbox{(§ 9)}
\]
이 따른다. \(\ml_i\mr_i\)는 0이 아닌 양쪽아이디얼이다(\(e_i^2\)가 \(\ml_i\mr_i\) 안에 있기 때문이다). 따라서 \(\ml_i\mr_i=\mo\)이다.

\(\mA_{ik}=\mr_i\ml_k\)로 놓으면
\[
        \mo=(\mr_1,\ldots,\mr_n)\cdot(\ml_1,\ldots,\ml_n)
        =(\ldots,\mA_{ij},\ldots)
\]
은 직접합이다. 실제로 \(0=\sum_{i,j}a_{ij}\)이면 \(a_{ij}\in\mr_i\)이고 \(\sum\mr_i\)가 직접합이므로
\[
        0=\sum_j a_{ij},
\]
이며, 다시 \(a_{ij}\in\ml_j\)이고 \(\sum\ml_j\)가 직접합이므로
\[
        0=a_{ij}
\]
이다.

또한 \(\mA_{ij}\mr_j=\mr_i\ml_j\mr_j=\mr_i\mo=\mr_i\)이므로 \(\mA_{ij}\ne0\)이다. \(\mA_{ij}\)에서 \(a_{ij}\ne0\)을 택하자. 그러면
\[
        a_{ij}\mr_j\subseteq\mr_i,
        \qquad
        a_{ij}\mr_j\ne0\quad\hbox{(\(a_{ij}e_j=a_{ij}\)이므로)}
\]
이고, 따라서
\[
        a_{ij}\mr_j=\mr_i.
\]
각 \(r_j\in\mr_j\)에 대해
\[
        a_{ij}r_j=r_i
\]
로 놓으면 사상 \(r_j\mapsto r_i\)는 \(\mr_j\)에서 \(\mr_i\)로 가는 연산자준동형이다. 0에 대응되는 원소들은 \(\mr_j\)에 포함되는 아이디얼을 이루므로 영아이디얼이다. 따라서 이 사상은 \emph{동형}이다. 그러므로 다음을 얻는다.

\emph{\(\mo\)의 한 분해에 나타나는 임의의 두 단순 우측아이디얼 \(\mr_i,\mr_j\)는 연산자동형이며, 동형들은 \(\mA_{ij}\)의 원소들에 의해 주어진다.} 서로 다른 임의의 두 단순 우측아이디얼 \(\mr,\mr'\)은 언제나 \(\mo\)의 적어도 한 분해에 함께 나타나므로, 그러한 임의의 두 아이디얼도 동형이다.

\emph{\(\mr_j\)에서 \(\mr_i\)로 가는 모든 준동형은 \(\mA_{ij}\)의 원소들에 의해 주어진다.}

\emph{증명.} \(e_j\mapsto a_{ij}\)이면 \(e_j^2\mapsto a_{ij}e_j\)이므로 \(a_{ij}=a_{ij}e_j\)는 \(\mr_i\ml_j=\mA_{ij}\)에 속한다. 또한
\[
        \mr_j=e_j\mr_j\longmapsto a_{ij}\mr_j=\mr_i.
\]

\emph{특히 \(\mr_i\)에서 자기 자신으로 가는 준동형들은 \(\mA_{ii}\)의 원소들에 의해 주어진다.}

\(\mA_{ii}\)의 두 원소의 곱에는 준동형들의 곱이, 합에는 그 합이 대응한다(정의는 § 1, 6 참조). 서로 다른 \(a_{ii}\)는 서로 다른 자기동형을 준다. 실제로 \(e_i\)는 \(a_{ii}e_i=a_{ii}\)로 간다. 따라서 환 \(\mA_{ii}\)는 \(\mr_i\)의 자기동형환과 동형이다.

그런데 단순아이디얼의 자기동형환은 체이다(§ 2). 따라서 \(\mA_{ii}\)는 체이다. 또한 모든 \(\mr_i\)가 연산자동형이므로 그 자기동형환들도 환동형이다. 곧 \(\mA_{ii}\)는 환동형을 제외하면 \(\mo\)에 의해 유일하게 결정된다. 가능한 서로 다른 \(\mA_{ii}\)들은 단순 우측아이디얼의 추상적으로 정의된 자기동형환을 구체적으로 실현하는 서로 다른 방식일 뿐이다.

\emph{주정리. 단위원을 갖는 모든 양쪽단순 완전가약환 \(\mo\)는 어떤 체 \(K\)의 원소로 된 \(n\)행 정방행렬환과 동형이다. (체 \(K\)는 \(\mo\)의 우측아이디얼들의 자기동형체와 동형이다.)}

\emph{증명\textsuperscript{13a)}. 행렬단위 \(c_{ik}\)의 구성.} \(\Gamma_{11}\)을 \(\mr_1\)의 항등자기동형이라 하고, \(\Gamma_{i1}\mr_1=\mr_i\)인 임의의 동형 \(\Gamma_{i1}\)들을 택한 뒤
\[
        \Gamma_{ik}=\Gamma_{i1}\Gamma_{k1}^{-1}
\]
로 놓자. 그러면 일반적으로
\[
        \Gamma_{ik}\Gamma_{kl}=\Gamma_{il}
\]
이다. \(\Gamma_{ik}\)가 \(\mA_{ik}\)의 원소 \(c_{ik}\)에 의해 주어지면
\[
        c_{ik}=e_i c_{ik}=c_{ik}e_k,
        \qquad
        c_{ik}c_{kl}e_l=c_{il}e_l,
\]
따라서
\[
        c_{ik}c_{kl}=c_{il},
        \qquad
        c_{ij}c_{kl}=c_{ij}e_j e_k c_{kl}=0\quad(j\ne k)
\]
이다. 그러므로 \(c_{ik}\)들은 행렬단위이다(§ 7).

\footnotetext{최근 B. L. van der Waerden은 추상적 자기동형체를 직접 써서 더 간단히 증명하는 방법을 찾아냈다. 이 증명은 그의 대수학 저서(\emph{Grundlehren der mathematischen Wissenschaften}, Berlin: Julius Springer)에 실릴 예정이다. [1929. 6. 11.]}

\emph{\(K\)의 구성.} 대응
\[
        a_{ii}=c_{i1}a_{11}c_{1i}
\]
는 \(\mA_{11}\)의 각 \(a_{11}\)에 \(\mA_{ii}\)의 \glqq{}켤레원소\grqq{} \(a_{ii}\)를 대응시킨다.

이 대응은 환동형이다. \(a\)와 \(c\)에 대응하는 동형 \(\Delta,\Gamma\)가
\[
        \Delta_{ii}=\Gamma_{i1}\Delta_{11}\Gamma_{i1}^{-1}
\]
로 연결되기 때문이다. 이 관계가 환동형을 나타낸다는 것은 잘 알려져 있다. 이제 원소
\[
        \alpha=a_{11}+\cdots+a_{nn}
        =a_{11}+c_{21}a_{11}c_{12}+\cdots+c_{n1}a_{11}c_{1n}
\]
를 만들자. 각 \(a_{11}\)에는 하나의 \(\alpha\)가 대응하며, 직접합이므로 이 대응은 일대일이다. 모든 \(\alpha\)의 전체는 체 \(K\sim\mA\)이다. 실제로
\[
        \alpha=a_{11}+\cdots+a_{nn},
        \qquad
        \beta=b_{11}+\cdots+b_{nn}
\]
이면
\[
        \alpha+\beta=(a_{11}+b_{11})+\cdots+(a_{nn}+b_{nn}),
        \qquad
        \alpha\beta=a_{11}b_{11}+\cdots+a_{nn}b_{nn}
\]
이다. 따라서 대응 \(a_{11}\mapsto\alpha\)는 동형이다.

또
\[
        c_{ik}\alpha=c_{ik}a_{kk}
        =c_{ik}c_{k1}a_{11}c_{1k}=c_{i1}a_{11}c_{1k},
\]
\[
        \alpha c_{ik}=a_{ii}c_{ik}
        =c_{i1}a_{11}c_{1i}c_{ik}=c_{i1}a_{11}c_{1k},
\]
이므로 각 \(\alpha\)는 모든 \(c_{ik}\)와 가환한다.

끝으로 같은 식들에서
\[
        c_{ik}K=c_{ik}\mA_{kk}
        =c_{ik}\mr_k\ml_k=\mr_i\ml_k=\mA_{ik}
\]
가 따르며, 따라서
\[
        \mo=\sum c_{ik}K.
\]

\emph{대응.} \(\mo\)의 각 원소를 \(\sum c_{ik}\alpha_{ik}\) 꼴로 나타내어 그 원소에 행렬 \((\alpha_{ik})\)를 대응시키면, 행렬곱의 정의에서 곧바로 알 수 있듯 이 대응은 동형이다. 이로써 주정리가 증명되었다.

\emph{역. 체 \(\Lambda\)의 원소로 된 \(n\)행 정방행렬환은 양쪽단순이고 완전가약이다. \(\Lambda\)는 우측아이디얼들의 자기동형환과 동형이고, \(\Lambda e\)는 앞 정리에서 구성한 \(K\)이다.}

\emph{증명.} 행렬단위(§ 7)를 \(c_{ik}\)라 하고
\[
        \mr_i=\sum_k c_{ik}\Lambda
\]
로 놓자. 그러면 \(\mo=\mr_1+\cdots+\mr_n\)임은 명백하다. \(\mr_i\)들은 단순 우측아이디얼이다. 실제로 0이 아닌 모든 원소 \(\sum c_{ik}\alpha_k\)는 \(\mr_i\) 전체를 생성한다. \(\alpha_j\ne0\)이면
\[
        \Bigl(\sum c_{ik}\alpha_k\Bigr)
        (c_{jl}\alpha_j^{-1})=c_{il}
\]
이고, \(\sum c_{il}\beta_l\)은 \(\mr_i\)의 모든 원소를 지난다.

\(\mr_i\)들은 연산자동형이다. 곧 \(\mr_i=c_{ik}\mr_k\)이다.

여기서 \(\mo\)가 양쪽으로 분해불능임이 따른다. 따라서 이미 증명한 완전가약성과 단위원의 존재에 의해 § 13의 정리로부터 \(\mo\)는 양쪽단순이다. 이는 0이 아닌 \(\sum\sum c_{ik}\alpha_{ik}\) 하나가 양쪽아이디얼 \(\mo\) 전체를 생성한다는 사실에서도 알 수 있다.

또한
\[
        \mr_i=c_{ii}\mo,
        \qquad
        \ml_i=\mo c_{ii},
\]
\[
        e=\sum c_{ii},
        \qquad c_{ii}=e_i,
\]
\[
        \mA_{ii}=\mr_i\cap\ml_i=c_{ii}\Lambda\sim\Lambda.
\]
끝으로 \(a_{11}=c_{11}\lambda\)로 놓으면
\[
        \alpha=a_{11}+c_{21}a_{11}c_{12}+\cdots
        =c_{11}\lambda+c_{22}\lambda+\cdots=e\lambda
\]
이므로 모든 것이 증명되었다.

\emph{\(\mo\)의 중심은 \(K\)의 중심이며, 따라서 가환체이다.}

\emph{증명.} \(\mZ(K)\)는 모든 \(x\)와 가환하는 \(\zeta\)들의 전체이다. 그런데 이들은 모든 \(c_{ik}\)와도 가환하며, 따라서 모든 합 \(\sum c_{ik}a_{ik}\)와도 가환한다. 그러므로
\[
        \mZ(K)\subseteq\mZ(\mo).
\]
역으로 \(z\in\mZ(\mo)\)이고
\[
        z=\sum_\mu\sum_\nu c_{\mu\nu}\gamma_{\mu\nu}
\]
라 하자.
\[
        zc_{ik}=c_{ik}z,
        \qquad
        \sum_\mu c_{\mu k}\gamma_{\mu i}
        =\sum_\nu c_{i\nu}\gamma_{k\nu}.
\]
따라서
\[
        \gamma_{\mu i}=0\quad(\mu\ne i),
        \qquad
        c_{ik}\gamma_{ii}=c_{ik}\gamma_{kk},
        \qquad
        \gamma_{ii}=\gamma_{kk}=\gamma,
\]
\[
        z=\sum c_{ii}\gamma=e\gamma=\gamma.
\]
따라서 \(z\in K\)이고 모든 \(x\)와 가환하므로 \(z\in\mZ(K)\)이다.

행렬환은 물론 좌측 완전가약이기도 하고, 단위원을 갖는 모든 우측 완전가약환은 행렬환들의 직접합이므로 다음이 따른다.

\emph{단위원을 갖는 모든 우측 완전가약환은 좌측으로도 완전가약이며, 그 역도 성립한다.}

또 다음이 성립한다. \emph{단위원을 갖는 완전가약환의 중심은 양쪽단순 행렬환들에 대응하는 가환체들의 직접합이다.}

다음 정리도 성립한다.\textsuperscript{14)}

\emph{서로 다른 두 분해 \(\mo=\sum c_{ik}K\), \(\mo=\sum c'_{ik}K'\)는 내부자기동형 \(\alpha'\mapsto x^{-1}\alpha x\), \(c'_{ik}=x^{-1}c_{ik}x\)에 의해 서로 옮겨진다.}
\footnotetext{E. Artin, 앞의 논문, 258쪽 참조.}

\emph{증명.} \(\mo=\mr_1+\cdots+\mr_n=\mr_1'+\cdots+\mr_n'\)이라 하고, \(a,b\)를 아이디얼 \(\mr_1,\mr_1'\) 사이의 서로 역인 두 동형을 주는 원소들이라 하자.
\[
        \mr_1=a\mr_1',
        \qquad
        \mr_1'=b\mr_1,
        \qquad
        ab=c_{11},
        \qquad
        ba=c'_{11}.
\]
\[
        x=\sum_i c_{i1}ac'_{1i},
        \qquad
        y=\sum_i c'_{i1}bc_{1i}
\]
로 놓으면
\[
        xy=e,
        \qquad
        y=x^{-1},
        \qquad
        x^{-1}c_{ik}x=c'_{ik},
        \qquad
        x^{-1}\alpha x=\alpha'
\]
이다.

\section*{제III장. 가군론과 표현론.}

\subsection*{§ 15. 표현과 표현가군\textsuperscript{15)}.}

\(\mo\)를 환, \(K\)를 단위원을 갖는 환이라 하자. (뒤의 응용에서는 \(K\)가 언제나 체이다.)

\emph{\(\mo\)의 \(K\) 안에서의 \(n\)차 표현이란 준동형}
\[
        \mo\sim\mathfrak D
\]
\emph{을 말한다. 여기서 \(\mathfrak D\)는 \(K\)의 원소로 된 \(n\)행 정방행렬들이 이루는 환이다.}

\(\mo\)의 \(K\)에 관한 \emph{표현가군}이란 좌측 \(\mo\)-가군이자 우측 \(K\)-가군인 쌍가군 \(\mM\)(§ 1, 5)을 말한다.
\[
        \mo\mM\subseteq\mM,
        \qquad
        \mM K\subseteq\mM.
\]
또한 이는 유한개의 한 원소 \(K\)-가군의 직접합
\[
        \mM=x_1K+\cdots+x_nK
\]
이어야 하며, \(K\)의 단위원은 항등연산자이어야 한다.

\emph{모든 표현가군은 하나의 표현을 낳는다.} \(c\in\mo\)이고
\[
        cx_k=\sum_i x_i\gamma_{ik},
\]
또는
\[
        c(x_1,\ldots,x_n)=(x_1,\ldots,x_n)C,
        \qquad C=(\gamma_{ik})
\]
라 하자. 그러면 행렬 \(C\)들은 \(c\)의 표현을 이룬다. 실제로
\[
        (b+c)x_k=\sum_i x_i(\beta_{ik}+\gamma_{ik}),
\]
\[
\begin{aligned}
        bcx_k
        &=b\sum_jx_j\gamma_{jk}
          =\sum_j bx_j\gamma_{jk}
          =\sum_j\sum_i x_i\beta_{ij}\gamma_{jk}\\
        &=\sum_i x_i\biggl(\sum_j\beta_{ij}\gamma_{jk}\biggr),
\end{aligned}
\]
또는
\[
        bc(x_1,\ldots,x_n)=(x_1,\ldots,x_n)BC
\]
이기 때문이다.

\emph{역으로, \(\mo\)의 모든 표현 \(\mathfrak D\)에는 하나의 표현가군, 더 정확히는 그 가군의 한 특정 기저가 속한다.} 실제로 \(\mM\)을 \(x_1,\ldots,x_n\)의 형식적 선형형식
\[
        y=\sum_i x_i\alpha_i
\]
전체로 이해하자. 그러면 \(\mM\)은 우측 \(K\)-가군이다. 원소 \(c\)에 행렬 \(C=(\gamma_{ik})\)가 대응할 때, 더 나아가
\srcnumdisplay[3.2em]{(1)}{%
\begin{gathered}
        cx_k=\sum_i x_i\gamma_{ik},\\
        c\biggl(\sum_k x_k\alpha_k\biggr)
        =\sum_i\sum_k x_i\gamma_{ik}\alpha_k
\end{gathered}}
로 정의하자. 준동형관계
\[
        c+d\longmapsto C+D,
        \qquad cd\longmapsto CD
\]
에서
\[
        (c+d)x_i=cx_i+dx_i,
        \qquad cdx_i=c\cdot dx_i
\]
가 따르고, 합 \(\sum_i x_i\alpha_i\)에 대해서도 마찬가지이다. 나머지 쌍가군 성질
\[
        c(y+z)=cy+cz,
        \qquad c(y\alpha)=(cy)\alpha
\]
은 자명하다. 따라서 \(\mM\)은 표현가군이고, (1)에 의해 정확히 표현 \(\mo\sim\mathfrak D\)에 속한다.

이제 \((y_1,\ldots,y_n)\)을 같은 표현가군의 다른 기저라 하자.
\[
        (y_1,\ldots,y_n)=(x_1,\ldots,x_n)P,
        \qquad
        (x_1,\ldots,x_n)=(y_1,\ldots,y_n)P^{-1}.
\]
\(x\)라는 기저에 관하여 \(b\sim B\)이면
\[
        b(y_1,\ldots,y_n)=b(x_1,\ldots,x_n)P
        =(x_1,\ldots,x_n)BP
        =(y_1,\ldots,y_n)P^{-1}BP.
\]
표현 \(b\sim B\)와 \(b\sim P^{-1}BP\)를 \emph{동치}라 하고 같은 \emph{표현류}에 넣는다. 모든 정칙행렬 \(P\)는 \(\mM\)의 기저를 다시 기저로 옮기므로 다음이 증명되었다.

\emph{모든 표현가군은 하나의 특정 표현류를 유일하게 결정한다.}\textsuperscript{15a)}

{\footnotesize\noindent
15) 가환체에 관한 표현가군은 W. Krull, \emph{Theorie und Anwendung der verallgemeinerten Abelschen Gruppen}, Sitzungsber. Heidelberger Ak. 1926, 제1논문에서 처음 나타난다.\par
15a) 교정 때의 추가(1929년 4월 14일). B. L. van der Waerden이 내게 알려 준 대로, \emph{선형변환과 행렬}이라는 개념을 분리하면 특정한 기저 선택과 무관한 불변적 관계를 직접 얻을 수 있다. 선형변환은 두 선형형식가군 사이의 준동형이고, 행렬은 특정한 기저 선택에서 이 준동형을 나타낸 표현이다. \emph{I. 모든 표현가군은 승수영역 \(\mo\)에 그 가군에서 자기 자신으로 가는 선형변환들의 한 준동형계를 유일하게 대응시킨다.} 실제로 § 2 끝부분에 따라 쌍가군 \(\mM\)의 좌측승수들은 우측연산자(여기서는 \textbf{\(K\)와의 곱셈})에 관하여 \(\mM\)에서 자기 자신으로 가는 연산자준동형을 생성하며, 그 대응은 준동형이다. \emph{II. \(\mo\)와 준동형인, 한 선형형식가군에서 자기 자신으로 가는 모든 선형변환계는 표현가군을 낳는다.} 다시 § 2 끝부분에 따라 이 준동형들은 좌측승수로서 쌍가군을 낳는다. 따라서 모든 \(m\in\mM\), \(c\in\mo\)에 대해 \(cm=\Gamma m\)로 놓으면 승수영역 \(\mo\)에 관해서도 마찬가지이다. 여기서 \(c\mapsto\Gamma\)는 처음 주어진 환준동형 대응이다. 우측ㆍ좌측 곱셈에 관하여 연산자동형인 표현가군들에는 같은 변환들이, 연산자동형이 아닌 표현가군들에는 서로 다른 변환들이 대응한다. 특정한 기저를 골라 각 변환을 행렬로 나타내면 I은 표현 \(\mo\sim\mathfrak D\)를 주고, II는 그러한 모든 표현이 하나의 표현가군, 더 정확히는 그 가군의 한 특정 기저를 낳는다고 말한다. 여기서 표현가군과 표현류의 대응이 다시 따르며, 계산은 전혀 필요하지 않다.\par}

다음은 명백하다.

\emph{연산자동형인 두 표현가군에는 같은 표현류가 대응한다.}

그러나 역도 성립한다.

\emph{두 표현가군 \((x_1,\ldots,x_n)\)과 \((\bar x_1,\ldots,\bar x_n)\)이 같은 표현을 생성하면, 대응 \(x_i\mapsto\bar x_i\), \(\sum_i x_i\alpha_i\mapsto\sum_i\bar x_i\alpha_i\)는 동형을 준다.} 곱셈법칙이 행렬 \(C\)들에 의해 완전히 결정되기 때문이다.

특수한 경우로, \(\mM\)의 두 기저가 같은 표현을 생성하면 이들은 \(\mM\)에서 자기 자신으로 가는 한 연산자동형에 의해 서로 옮겨진다.

달리 말해, 고정된 \(P\)에 대해 모든 \(C\)가 \(C=PCP^{-1}\)을 만족하면 \(y_i=\sum_kx_k\pi_{ik}\)일 때 대응 \(y_i\mapsto x_i\)는 \(\mM\)에서 자기 자신으로 가는 동형이다. 역으로 쌍가군 \(\mM\)에서 자기 자신으로 가는 모든 동형 \(y_i\mapsto x_i\)는 모든 \(C\)와 가환하는 행렬 \(P\)에 의해 주어진다.

표현 개념은 가환 승수영역 \(P\)와 가환적으로 결합된 환들로도 확장할 수 있다(§ 8). 이 경우에는 중심에 \(P\)를 포함하는 표현환 \(K\)만 취하고, 표현이 환준동형 \(\mo\sim\mathfrak D\)일 뿐 아니라 연산자준동형일 것도 요구한다. 즉 \(a\mapsto A\)이면 모든 \(\rho\in P\)에 대해 \(a\rho\mapsto A\rho\)이어야 한다.
표현가군에 관해서 이 요구는 계산법칙
\[
        am\cdot\rho\ [=a\cdot m\rho]=a\rho\cdot m
\]
을 추가한다. 이때 가군과 환은 흔히 \glqq{}\(P\)와 가환적으로 결합되어 있다\grqq{}고 한다.

\subsection*{§ 16. 가약표현.}

\emph{\(\mA\)가 표현가군 \(\mM\)의 부분가군이고, \(\mA\)의 기저 \(z_1,\ldots,z_t\)에 \(y_1,\ldots,y_r\)을 덧붙여 \(\mM\)의 기저를 고를 수 있어서}
\[
        \mM=y_1K+\cdots+y_rK+z_1K+\cdots+z_tK\textsuperscript{16)}
\]
\emph{이면, 표현들은}
\[
        C=\begin{pmatrix}R&0\\ S&T\end{pmatrix}
\]
\emph{의 꼴이다. 여기서 행렬 \(T\)들만으로 \(\mA\)가 생성하는 \(t\)차 표현을 이루고, 행렬 \(R\)들은 \(\mM/\mA\)가 생성하는 \(r\)차 표현을 이룬다.}

\emph{증명.} 
\[
        (cy_1,\ldots,cy_r,cz_1,\ldots,cz_t)
        =(y_1,\ldots,y_r,z_1,\ldots,z_t)C.
\]
\(cz_i\)들은 \(\mA\)의 원소이므로 \(z_i\)들만으로 나타난다. 따라서 \(C\)의 오른쪽 위 블록은 0이다. \(C\)의 나머지 블록들을 위 배열과 같이 \(R,S,T\)라 하면, 특히
\[
        (cz_1,\ldots,cz_t)=(z_1,\ldots,z_t)T
\]
이다. 그러므로 \(T\)들은 \(\mA\)가 매개하는 표현을 이룬다. 또한
\[
        (cy_1,\ldots,cy_r)\equiv
        (y_1,\ldots,y_r)R\pmod{\mA},
\]
이고 \(y_i\)들은 \(\bmod\,\mA\) 선형독립기저를 이룬다. 따라서 \(R\)들은 \(\mM/\mA\)가 생성하는 표현을 이룬다.

역으로
\[
        C=\begin{pmatrix}R&0\\ S&T\end{pmatrix}
\]
꼴의 \glqq{}가약\grqq{} 표현이 주어지고 \(R,T\)가 정방행렬이면, 이에 속하는 표현가군에서 각 \(c\)와 마지막 \(t\)개 기저원소 \(z_1,\ldots,z_t\)의 곱은 이 \(z_i\)들만으로 나타난다. 곧 \(\mA=(z_1,\ldots,z_t)\)는 부분가군이다.

{\footnotesize\noindent 16) 달리 말하면 \(\mA\)가 \(K\)-가군으로서 직접가수라는 뜻이다. \(K\)가 체이면 이 가정은 언제나 성립한다.\par}

\emph{따름정리.}
\[
        \mM>\mA_1>\mA_2>\cdots>\mA_e=0
\]
을 \(\mM\)의 한 합성열이라 하고, 각 \(\mA_i\)가 앞 가군의 \(K\)-가군으로서 직접가수라 하자. 그러면
\[
        z_{11},\ldots,z_{1r_1},\ z_{21},\ldots,z_{2r_2},\ldots,
        z_{e1},\ldots,z_{er_e}
\]
라는 기저를 골라 \(i>\nu\)인 \(z_{ik}\)들이 \(\mA_\nu\)의 기저를 이루게 할 수 있다. 이때 \(\mM\)이 매개하는 표현은
\srcnumdisplay[3.2em]{(2)}{%
 C=\begin{pmatrix}
 R_{11}&0&\cdots&0\\
 R_{12}&R_{22}&\cdots&0\\
 \vdots&\vdots&\ddots&\vdots\\
 R_{1e}&R_{2e}&\cdots&R_{ee}
 \end{pmatrix}}
의 꼴이고, \(R_{\nu\nu}\)들은 합성인자 \(\mA_{\nu-1}/\mA_\nu\)가 매개하는 표현을 이룬다.

각 표현 \(R_{\nu\nu}\)은 기약이다. 합성인자 \(\mA_{\nu-1}/\mA_\nu\)가 단순하기 때문이다. \(K\)가 체라서 모든 부분가군이 직접가수라고 가정하면, 역으로 (2)의 꼴로 가능한 데까지 가약화된 모든 표현은 합성열을 낳는다. 조르당--횔더 정리에 따라 합성인자들은 연산자동형과 순서를 제외하면 유일하므로, \emph{\(R_{\nu\nu}\)들은 동치표현과 순서를 제외하면 유일하게 결정된다.}

\(\mM=\mA+\mB\)가 두 표현가군 \(\mA=(y_1,\ldots,y_r)\), \(\mB=(z_1,\ldots,z_s)\)의 직접합이면, 기저 \((y_1,\ldots,y_r,z_1,\ldots,z_s)\)가 매개하는 표현은 명백히
\[
        C=\begin{pmatrix}R&0\\0&S\end{pmatrix}
\]
의 꼴이다. 여기서 \(R\)은 \(\mA\)가, \(S\)는 \(\mB\)가 매개하는 원소 \(c\)의 표현이며, 역도 성립한다. 먼저 가군 \(\mM\)을 직접분해불능 성분들의 직접합으로 나타내고, 각 성분에서 위와 같은 합성열을 찾으면 알맞은 기저를 골랐을 때 표현은
\[
\begin{array}{@{}r@{\quad}c@{}}
(3)&
C=\left(
\begin{array}{ccc}
\boxed{\begin{matrix}
R_{11}&0\\
\vdots&\ddots\\
R_{ik}&R_{rr}
\end{matrix}} & & 0\\[1.4em]
& \boxed{\begin{matrix}
S_{11}&0\\
\vdots&\ddots\\
S_{ik}&S_{ss}
\end{matrix}} & \\[1.4em]
0 & & \boxed{\begin{matrix}
T_{11}&0\\
\vdots&\ddots\\
T_{ik}&T_{tt}
\end{matrix}}
\end{array}
\right)
\end{array}
\]
의 꼴이다.

다시 \(K\)가 체이면, 역으로 (3)의 꼴로 가능한 데까지 가약화된 모든 표현은 가군을 직접분해불능 가수들로 나타내고 그 안에서 합성인자들을 택하는 분해를 낳는다. 크룰과 오토 슈미트의 정리(각주 11)에 따라 직접분해불능 표현성분들의 류는 순서를 제외하면 유일하게 결정된다.

가군 \(\mM\)이 완전가약이면 모든 상자는 하나의 기약표현만으로 이루어지며, 표현 (3)을 완전가약이라 한다.

\subsection*{§ 17. 단위원을 갖는 환에서 표현가군의 직접합분해.}

\emph{\(\mM\)을 \(\mo\)-가군, \(\mo\)를 단위원을 갖는 환이라 하자. 그러면}
\[
        \mM=\mM_e+\mM_0,
\]
\emph{여기서 단위원은 \(\mM_e\)에서는 항등연산자이고 \(\mM_0\)에서는 영연산자이다. (\(\mM\)이 우측 \(K\)-가군이면 \(\mM_e,\mM_0\)도 그러하다.)}

\emph{증명.} \(\mM\)의 모든 \(m\)은
\[
        m=em+(m-em)
\]
으로 나타난다. \(em\) 전체를 \(\mM_e\), \(m-em\) 전체를 \(\mM_0\)라 하자. 그러면 \(\mM_e,\mM_0\)가 가법군이며, \(\mM\)이 우측 \(K\)-가군일 때 이들도 우측 \(K\)-가군이라는 것은 자명하다. 또한
\[
        r em=erm
\]
이므로 \(\mM_e\)도 \(\mo\)-가군이다. 마찬가지로
\[
        r(m-em)=rm-rem=rm-rm
\]
이므로 \(\mM_0\)도 \(\mo\)-가군이다. \(em\)들에는 \(e\)가 항등연산자이고 \(m-em\)들에는 영연산자이다. 따라서 \(\mM_e\)와 \(\mM_0\) 양쪽에 속하는 것은 0뿐이며 합 \(\mM_e+\mM_0\)는 직접합이다.

표현가군의 경우 \(\mM_0\)는 모든 원소에 0을 대응시키는 자명표현을 생성한다. 가수 \(\mM_0\)를 떼어 내면 \(\mM_e\)가 매개하는 표현이 남고, \emph{이 표현에서는 단위원에 단위행렬이 대응한다.} 그러므로 이제부터는 그러한 표현에만 국한할 수 있고 또 그렇게 하겠다.

\emph{\(\mo=\ma_1+\cdots+\ma_s\)가 양쪽아이디얼들의 직접합이고, \(e=e_1+\cdots+e_s\)가 \(e\)의 분해이며, \(\mM\)이 \(e\)를 항등연산자로 갖는 \(\mo\)-가군이면}
\[
        \mM=\ma_1\mM+\cdots+\ma_s\mM
\]
\emph{은 직접합이다.} 실제로 \(\mM=\mo\mM=(\ma_1\mM,\ldots,\ma_s\mM)\)임은 명백하다. \(\mb_i=\ma_1+\cdots+\widehat{\ma_i}+\cdots+\ma_s\)로 놓으면
\[
        \mM=(\ma_i\mM,\mb_i\mM)
\]
이고, \(e_i\)가 \(\ma_i\mM\)에서는 항등연산자, \(\mb_i\mM\)에서는 영연산자이므로 이 합은 직접합이다.

이 정리에 따라 우리는 대개 양쪽으로 분해불능인 환의 표현에 국한한다.

\subsection*{§ 18. 완전가약환의 가군론과 표현론.}

\emph{\(\mo\)가 완전가약이고 양쪽단순이며, \(\mM\)이 \(\mo\)의 단위원을 항등연산자로 갖는 유한 \(\mo\)-가군이라고 하자. 그러면 \(\mM\)은 완전가약이고, 그 단순성분들은 단순 좌측아이디얼 \(\ml_i\)들과 연산자동형이다.}

\emph{증명.} 
\[
        \mo=\ml_1+\cdots+\ml_n,
        \qquad
        \mM=(\mo m_1,\ldots,\mo m_k)
\]
라 하자. 그러면
\[
\makebox[0pt][r]{(4)\qquad}\mM=(\ldots,\ml_i m_k,\ldots).
\]
0이 아닌 \(\ml_i m_k\)들은 \(\ml_i\)와 동형이다. 대응 \(a\mapsto am_k\)가 연산자동형이기 때문이다. 따라서 이 \(\ml_i m_k\)들은 단순하다. 표시 (4)에서 앞선 항들의 합에 이미 포함되는 \(\ml_i m_k\)들을 버리면 합은 직접합이 된다.

\emph{따름정리.} 또한 \(\mM\)이 직접분해불능이면 \(\mM\)은 단순이고 어떤 \(\ml_i\)와 동형이다.

이 단순 \(\mM\)의 자기동형체를 \(\Delta\), \(\ml_i\)의 자기동형체를 \(\Lambda\)라 하자. \(\mM\)과 \(\ml_i\) 사이의 연산자동형에 의해 환동형 \(\Delta\simeq\Lambda\)가 성립한다. § 1에 따라 \(\mM\)과 \(\ml_i\)를 각각 \(\Delta\), \(\Lambda\)를 우측영역으로 하는 쌍가군으로 볼 수 있으며, 이 쌍가군들은 표현가군이기도 하다. 실제로 \(\ml_i\), 따라서 \(\mM\)은 자기동형체에 관하여 유한랭크이고(§ 14), 그 체의 단위원은 항등연산자이다. 이들이 매개하는 표현들은 환동형 \(\Delta\simeq\Lambda\)에 의해 서로 옮겨진다. \emph{따라서 이 표현류를 연구할 때에는 \(\ml_i\)가 그 자기동형체에서 매개하는 표현에 국한할 수 있다.}

특히 \(\ml_1\)에 대해 행렬단위들로 된 기저
\[
        \ml_1=(c_{11},\ldots,c_{n1})
\]
를 택하고, 자기동형체 \(\Lambda\)의 실현으로 앞에서 \(K\)라 부른 \(\mo\)의 부분체(§ 14)를 택하자. 이때 앞의 논의에서 우측아이디얼을 좌측아이디얼로 바꾸고, 그에 따라 자기동형들을 오른쪽에 써야 한다. 
\[
        a=\sum c_{ik}\alpha_{ik}
\]
로 나타내면
\[
        a\longmapsto(\alpha_{ik})
\]
는 \(\ml_1\)이 매개하는 \(K\) 안에서의 표현이다. 실제로
\[
\begin{aligned}
        a\cdot c_{k1}
        &=\left(\sum\sum c_{ij}\alpha_{ij}\right)c_{k1}\\
        &=\sum c_{ik}\alpha_{ik}c_{k1}
         =\sum c_{i1}\alpha_{ik},
\end{aligned}
\]
또는
\[
        a(c_{11},\ldots,c_{n1})
        =(c_{11},\ldots,c_{n1})(\alpha_{ik})
\]
이다.

\emph{\(\Gamma\)를 \(K\)의 부분체라 하고, \(K\)가 \(\Gamma\)에 관하여 유한랭크 \(t\)를 갖는다고 하자.}
\[
        K=z_1\Gamma+\cdots+z_t\Gamma.
\]
\emph{그러면 \(K\)는 \(\Gamma\)에 관한 자기 자신의 표현가군이 된다. \(K\)의 원소 \(\beta\)는}
\[
        \beta z_j=\sum_i z_i\beta_{ij},
        \qquad B=(\beta_{ij})
\]
\emph{로 얻는 \(\Gamma\)의 행렬 \(B\)에 의해 표현된다.}

이제 \(\ml_1\)을 부분체 \(\Gamma\)에 관한 표현가군으로 보면 다음을 얻는다.

\emph{\(\ml_1\)이 매개하는 \(\mo\)의 \(\Gamma\) 안에서의 표현은, 위에서 준 \(\mo\)의 \(K\) 안에서의 표현}
\[
        a\longmapsto(\alpha_{ik})
\]
\emph{에서 원소 \(\alpha_{ik}\)를 \(K\)의 \(\Gamma\) 안에서의 표현에서 이에 대응하는 행렬 \(A_{ik}\)로 바꾸어 얻는다.}

\emph{증명.}
\[
        \ml_1=\sum c_{i1}K=\sum\sum c_{i1}z_j\Gamma.
\]
기저 \((\ldots,c_{i1}z_j,\ldots)\)에 관하여
\[
c_{ik}\longmapsto
\begin{array}{c@{\quad}c}
 & \begin{array}{ccc}
 \boxed{\scriptstyle 1}&\boxed{\scriptstyle k}&\boxed{\scriptstyle n}
 \end{array}\\[-.2em]
\begin{array}{c}
\boxed{\scriptstyle 1}\\ \boxed{\scriptstyle i}\\ \boxed{\scriptstyle n}
\end{array}
&
\begin{pmatrix}
0&\cdots&0&\cdots&0\\
\vdots&&\vdots&&\vdots\\
0&\cdots&E_t&\cdots&0\\
\vdots&&\vdots&&\vdots\\
0&\cdots&0&\cdots&0
\end{pmatrix}
\end{array},
\]
이다. 여기서 \(E_t\)는 단위행렬이고 0은 \(t\)차 영행렬이다. 또한
\[
        \alpha\longmapsto
        \begin{pmatrix}
        A& &0\\
         &\ddots&\\
        0& &A
        \end{pmatrix}.
\]
따라서
\[
        a=\sum c_{ik}\alpha_{ik}\longmapsto
        \begin{pmatrix}
        A_{11}&\cdots&A_{1n}\\
        \vdots&&\vdots\\
        A_{n1}&\cdots&A_{nn}
        \end{pmatrix}.
\]

임의의 \(\mM\)으로 넘어가는 것은 다시 표현환에 환동형을 가하는 것을 뜻한다. \(\mo\)가 양쪽단순환들의 직접합이면 이 환들은 차례로 동형으로 표현되고, 나머지 환들은 0으로 표현된다.

\textbf{논평.} 여기서 연구한 좌측아이디얼의 자기동형체와 그 유한부분체 안에서의 표현들은, 단순 \(\mo\)-가군(또는 좌측아이디얼)이 매개할 수 있는 표현들 가운데 동형을 제외하면 유일한 것들이다. 실제로 앞에서 한 논평(§ 2)에 따라, \(\mo\)-가군 \(\mM\)을 임의의 체 \(K\)에 관한 쌍가군으로 보면 \(K\)의 원소들은 가군준동형들을 생성한다. 따라서 체 \(K\)는 이 \(\mo\)-가군의 자기동형체의 한 부분체에 준동형으로 대응해야 한다. 단위원이 항등연산자이므로 이 준동형은 동형이다. 또한 전체 자기동형체는 해당 부분체에 관하여 유한차수이어야 한다. 그렇지 않으면 표현가군도 무한랭크가 되어 불가능하기 때문이다.

끝으로 여기서 조사한 기약표현들이 전부는 아니고, 단순 \(\mo\)-가군들이 매개하는 것들뿐이라는 점을 덧붙인다. 일반적으로는 표현가군들, 예컨대 부분체 \(\Sigma\)를 좌측영역으로 하고 자기 자신을 우측영역으로 하는 쌍가군으로 본 체 \(\Omega\)가 더 있다. 이들은 \(\mo\)-가군으로서는 가약이지만 쌍가군으로서는 기약이고, 그러므로 기약표현을 낳는다.

\subsection*{§ 19. 가군과 표현가군의 단순 합성인자.}

\(\mo\)를 좌측아이디얼에 대한 극대조건과 극소조건을 만족하는 환, \(\mc\)를 극대 멱영아이디얼이라 하자. 그러면 다음이 성립한다.

\emph{모든 단순 \(\mo\)-가군 \(\mM\)은 \(\mo\)에 의해 소멸되거나, 근기 없는 환 \(\mo/\mc\)의 한 단순 좌측아이디얼 \(\ml\)과 동형이다. 뒤의 경우 그 가군은 \(\mo/\mc\)의 양쪽아이디얼 가운데 \(\ml\)을 포함하지 않는 모든 것에 의해 소멸되고, 절대승수영역(§ 1)은 \(\ml\)을 포함하는 양쪽단순환과 동형이다.}

\textbf{증명.} \(\mc^e=0\)이라 하자. \(\mc\mM=0\)이어야 한다. 그렇지 않으면 \(\mc\mM=\mM\)이고, 따라서 \(\mM=\mc\mM=\mc^2\mM=\cdots=0\)이 되어 버리기 때문이다. 그러므로 \(\mM\)을 \(\mo/\mc\)-가군으로도 볼 수 있다. \(\mc\)에 따른 같은 잉여류의 모든 원소는 \(\mM\)의 원소를 곱할 때 같은 결과를 준다. 근기 없는 환인 \(\mo/\mc\)는 단위원을 가지므로, \(\mM\)은 \(\mo/\mc\)에 의해 소멸되는 가군과 그 단위원이 항등연산자인 가군의 직접합이다(§ 17). \(\mM\)이 단순이므로 이 두 가수 가운데 하나만 나타날 수 있다. 단위원이 항등연산자이면 § 18에서와 같이 단순 좌측아이디얼과의 동형이 따른다. 실제로 \(m\ne0\)이 \(\mM\)의 원소이면 \((\mo/\mc)m\ne0\)이고, 따라서 \(\mo/\mc\)의 적어도 한 단순 좌측아이디얼 \(\ml\)에 대해 \(\ml m\ne0\)이다. 그러면 \(\ml m\)은 \(\mM\) 전체이어야 한다. 나머지 주장들은 그러한 좌측아이디얼에 대해서는 명백하며, 그와 동형인 모든 가군으로 곧바로 옮겨진다.

\textbf{따름정리.} \emph{\(\mM\)이 합성열을 갖는 \(\mo\)-가군이면, 그 단순 합성인자들은 \(\mo\)에 의해 소멸되거나 \(\mo/\mc\)의 단순 좌측아이디얼들과 동형이다.}

\(P\)가 \(\mo\)와 가환적으로 결합된 환이면 모든 결과가 그대로 성립한다. 다만 가군 대신 우측의 \(P\)와 좌측의 \(\mo\)에 관한 쌍가군으로서
\[
        am\cdot\rho=a\cdot m\rho=a\rho\cdot m
        \qquad\hbox{(§ 15)}
\]
을 만족하는 것들을 고찰해야 한다. 또한 \(P\)의 단위원도 항등연산자이어야 한다. 부분가군으로는 \(P\)와의 곱셈을 허용하는 것들만을 언제나 고찰한다(아이디얼의 경우와 같다). 우리의 증명에서 구성한 부분가군 \(\mc\mM,\mc^2\mM,\ldots\)는 이 조건을 만족한다.

특히 \(P\)가 체이고 \(\mM\)이 유한랭크 \(P\)-가군이면 \(\mM\)은 동시에 표현가군이다. \(\mM\)이 매개하는 표현은 환준동형일 뿐 아니라 \(P\)에 관한 연산자준동형이다(§ 15). 이 성질을 갖는 \(P\) 안에서의 모든 표현은 그러한 가군들로 얻어진다. 기약표현에는 단순가군이 대응하며, 역도 성립한다. 따라서 이 절의 정리들은 곧바로 \(\mo\)의 \(P\) 안에서의 표현에 관한 정리로 해석할 수 있다. \emph{모든 기약표현은 영표현을 제외하면 \(\mo/\mc\)의 단순 좌측아이디얼들이 매개한다. 따라서 임의의 표현을 § 16에 따라 합성열로 가약화하면, 주대각선에는 0 이외에 \(\mo/\mc\)의 단순 좌측아이디얼들이 매개하는 표현들과 동치인 것들만 나타난다.}

\textbf{논평.} \(\mo\)와 \(P\)에 관하여 위의 가정들을 두면 \(\mo\) 자체에는 유한성 조건을 가정할 필요가 없다. 모든 표현은 동시에 절대승수영역의 동형표현이기 때문이다. 이 절대승수영역은 유한차 행렬환의 동형원상으로서 스스로 유한랭크 \(P\)-가군이 된다. 또 정의상 \(P\)-가군인 것들만 허용아이디얼로 고찰하므로, 이 절대승수영역에서는 극대조건과 극소조건이 성립한다.

\(P\)를 가환체라고 가정하면 이 절대승수영역은 \(P\)에 관한 초복소계가 된다. 주대각선에 0 이외의 것이 하나라도 나타나면 이 가정은 언제나 충족된다. 실제로 § 18의 논평에 따라 \(P\)는 \(\ml\)의 자기동형체의 한 부분체와 동형이고, 이를 \(\mo/\mc\)의 부분체 \(K\)로 실현할 때(§ 14) 가환적 결합 때문에 \(P\)는 \(K\)의 중심에 놓여야 한다.

\section*{제IV장. 군과 초복소계의 표현.}

\subsection*{§ 20. 초복소계의 분류.}

가환체 \(P\)에 관한 초복소계 \(\mo\)를 고찰한다. 약속에 따라 \(\mo\)의 아이디얼로는 \(P\)-가군만을 고찰하므로 좌측ㆍ우측아이디얼에 대한 극대조건과 극소조건이 성립한다(§ 8 참조). 따라서 제II장에서 전개한 환론 전체가 적용된다.

이하에서 초복소계 \(\mo\)의 \emph{표현}이라 하면 언제나 체 \(P\) 안에서의 표현으로서, \(a\mapsto A\)이면 모든 \(\rho\in P\)에 대해 \(a\rho\mapsto A\rho\)인 것, 곧 \(P\)에 관한 가군준동형을 뜻한다. 표현가군에 대해서는 이것이 \(a\rho\cdot m=a\cdot m\rho\)를 뜻한다(§ 15 끝 참조). 따라서 표현가군에는 § 19의 정리들이 적용되며, 그 결과 \emph{모든 기약표현은 \(\bar{\mo}=\mo/\mc\)의 단순 좌측아이디얼들이 매개하고, 따라서 서로 동치가 아닌 기약표현의 수는 분해 \(\bar{\mo}=\bar{\mo}e^{(1)}+\cdots+\bar{\mo}e^{(s)}\)의 양쪽단순 가수 \(\ma\)의 수와 같다.}

그러한 좌측아이디얼 \(\ml_\nu\)가 정의하는 \(\mo\)의 표현을 실제로 구하려면, 먼저 \(\mo\)의 원소들에서 \(\bar{\mo}\)의 대응 잉여류들로 넘어간다. \(P\)의 원소와의 곱셈은 \(\bar{\mo}\)의 모든 아이디얼에서 동형을 이루므로 체 \(P\)는 각 단순 좌측아이디얼의 자기동형체 \(K_\nu\)의 부분체 \(e^{(\nu)}P\)에 준동형이고, 단위원이 항등연산자이므로 사실 동형이다. 단순 좌측아이디얼 \(\ml_\nu\)가 매개하는 \(P\) 안에서의 표현은 먼저 이 부분체 \(e^{(\nu)}P\) 안에서 \(\ml_\nu\)가 매개하는 표현을 구한 뒤, 동형 \(e^{(\nu)}P\simeq P\)로 \(P\)에 옮겨 얻는다. 자기동형체 \(K_\nu\)는 \(P\)에 관하여 유한랭크이다. 따라서 이는 \(\ml_\nu\)의 자기동형체의 유한차수 부분체 안에서의 표현을 구하는 문제이다.

그러므로 구하는 표현을 얻기 위해 먼저 \(\bar{\mo}\)의 각 \(c\)를 양쪽성분들로
\[
        c=c_1+\cdots+c_s
\]
나타내자. 다른 \(c_i\)들은 아이디얼 \(\ml_\nu\)를 소멸시켜 영행렬로 표현되므로 \(c_\nu\)를 나타내는 행렬만 찾으면 된다. § 18에 따라 다음과 같이 구한다. \(c_\nu\)를 \(\ma_\nu\)의 행렬단위 \(c_{ik}^{(\nu)}\)로
\[
        c_\nu=\sum c_{ik}^{(\nu)}\alpha_{ik}^{(\nu)}
\]
나타내고, 행렬 \((\alpha_{ik}^{(\nu)})\)의 각 원소 \(\alpha_{ik}^{(\nu)}\)를 \(K_\nu\)가 매개하는 \(K_\nu\)의 \(P\) 안에서의 표현에서 이에 대응하는 행렬 \(A_{ik}^{(\nu)}\)로 바꾼다.

\textbf{논평.} 체 \(K_\nu\)는 \(P\)에 관하여 유한랭크이다. \(x\)의 거듭제곱들 사이에는 반드시 선형종속관계가 있으므로 각 원소 \(x\)는 \(e^{(\nu)}P\)의 계수를 갖는 방정식 \(f(x)=0\)을 만족한다. 특히 \(P\)가 대수적으로 닫혀 있으면 이 방정식은 일차인자들로 분해된다. \(K\)가 체이므로 \(x\)는 이미 한 일차인자의 영점이고, 따라서 \(x\in e^{(\nu)}P\)이다. 곧 \(K_\nu=e^{(\nu)}P\)이다. 이 경우에는 행렬 \((\alpha_{ik}^{(\nu)})\) 자체가 기약표현을 이룬다.

이 논평과 § 19 끝부분에서 \glqq{}번사이드 정리\grqq{}가 나온다.

\emph{\(P\)가 대수적으로 닫혀 있고 환 \(\mo\)와 가환적으로 결합되어 있다고 하자. 그러면 \(P\) 안에서의 모든 \(n\)차 기약표현 \(\mo\sim\mD\)는 정확히 \(n^2\)개의 선형독립 행렬을 포함한다.}\textsuperscript{17)}

이 정식화에서 통상적인 정식화가 곧바로 따른다. \glqq{}\(P\)의 원소로 된 \(n\)행 정방행렬들의 한 계가 임의의 두 원소와 함께 그 곱도 포함한다고 하자. 이 계는, 그 계의 모든 행렬의 원소 \(\alpha_{ik}\)들이 만족하는 \(P\) 계수의 선형동차관계
\[
        \sum \alpha_{ik}\rho_{ik}=0
\]
가 존재하지 않을 때 그리고 그때에만 기약이다.\grqq{} 실제로 그러한 행렬계에 모든 선형결합을 덧붙이면 환이자 \(P\)-가군을 얻으며, 이 환을 자기 자신의 표현으로 볼 수 있다.

\textbf{증명.} § 19에 따라 \(\mD\)는 단위원을 갖는 양쪽단순 완전가약환 \(\bar{\mo}\)와 동형이다. 그런데 그러한 환은 좌측아이디얼의 자기동형체에 관한 완전행렬환, 이를테면 모든 \(m\)행 정방행렬들의 환이고, 대수적 폐성 때문에 이 자기동형체는 \(P\)와 같다. \(\bar{\mo}\)의 모든 기약표현, 따라서 주어진 표현도 단순 좌측아이디얼이 생성한다. 그 좌측아이디얼의 랭크는 \(m\)이고 표현의 차수는 \(n\)이므로 \(m=n\)이다. 따라서 \(\bar{\mo}\), 그리고 \(\mD\)에는 정확히 \(n^2\)개의 선형독립 원소가 있다.

마찬가지로 \glqq{}일반화된 번사이드 정리\grqq{}\textsuperscript{18)}도 쉽게 따른다.

\emph{같은 가정 아래 \(\mD\)가 \(\mo\)의 완전가약표현이고, 차수가 \(n_1,n_2,\ldots,n_s\)인 서로 동치가 아닌 성분들로만 분해된다고 하자. 그러면 \(\mD\)는 \(P\)에 관하여 랭크 \(n_1^2+\cdots+n_s^2\)를 갖는다.}

\textbf{증명.} 다시 \(\bar{\mo}\)를 절대승수영역이라 하자. 이는 반드시 초복소계이다. 서로 동치가 아닌 표현들은 \(\bar{\mo}/\bar{\mc}\)의 서로 다른 양쪽단순 성분 \(\ma_i\)에 속하는 좌측아이디얼 \(\ml_1,\ldots,\ml_s\)이 생성한다. 여기서 \(\bar{\mc}\)는 근기이다. 따라서 주어진 표현은 \(\ml_1+\cdots+\ml_s\)가 생성하며, \(\ma_1+\cdots+\ma_s\)를 절대승수영역으로 갖는다. 위와 같은 랭크 고찰로부터 주장이 따른다.

이제 \(\mo\)를 임의의 체에 관한 근기 없는 초복소계라 하자. § 13에 따라 \(\mo\)는 완전가약이고 단위원을 갖는다. § 17과 § 18의 정리들에서 다음이 따른다. \emph{\(\mo\)의 모든 표현은 완전가약이다.}

역으로 다음이 성립한다. \emph{\(P\)와 가환적으로 결합된 환 \(\mo\)의 모든 완전가약표현은 근기 없는 초복소계를 절대승수영역으로 갖는다.}

\textbf{증명.} 절대승수영역 \(\bar{\mo}\)는 \(P\)의 행렬들로 된 환이자 \(P\)-가군과 동형이므로 초복소계이다. 근기 \(\bar{\mc}\)의 원소들은 § 19에 따라 모든 기약표현에서, 따라서 전체 표현에서 0으로 표현된다. 그러므로 \(\bar{\mc}=0\)이다.

(표현에 서로 동치인 성분들만 나타나고 이들이 모두 한 좌측아이디얼 \(\ml\)에 의해 매개되면 절대승수영역은 양쪽단순이다.)

초복소계 \(\mo\)의 \emph{정칙표현}이란 단위아이디얼 \(\mo\) 자체가 매개하는 표현을 말한다. (군환에서는 특히 군원소 \(u_1,\ldots,u_h\)로 된 기저를 택한다. § 6 참조.) \(\mo\)의 합성인자들 가운데 \(\mo/\mc\)의 모든 좌측아이디얼이 나타나므로, 정칙표현을 § 16의 식 (2)에 따라 합성열을 써서 알맞은 기저로 변환하면 모든 기약표현이 이미 그 대각행렬 \(R_{ii}\)들 안에 들어 있다.

기저원소 \(u_1,\ldots,u_h\)를 나타내는 행렬들을 알면 표현은 알려진다. 특히 군환에서 \(u_1,\ldots,u_h\)가 한 군의 원소들이면, 그 군의 모든 행렬 준동형표현은 군환의 표현도 생성한다. 따라서 군의 표현문제는 초복소계 표현문제의 특수한 경우이다.

{\footnotesize\noindent
17) 따라서 \(\mo\)를 초복소계라고 가정하지 않는다. 예컨대 무한군의 원소들을 매번 유한개씩 취한 모든 선형결합으로 이루어진 군환일 수 있다.\par
18) G. Frobenius 및 I. Schur, \emph{Über die Äquivalenz der Gruppen linearer Substitutionen}, Sitzungsber. Berlin 1906.\par}

\subsection*{§ 21. 바탕체의 확대. 중심의 표현.}

\(\mo=a_1P+\cdots+a_hP\)가 초복소계이고 \(\Omega\)가 바탕체 \(P\)의 확대체이면, 기저원소 \(a_i\)에 대한 예전 곱셈법칙을 그대로 써서
\[
        \mo\Omega=a_1\Omega+\cdots+a_h\Omega
\]
를 만들 수 있다.\textsuperscript{18a)} \(\mo\Omega\)는 다시 \(\Omega\)에 관한 초복소계이다.

\(\mo\)의 모든 표현은 \(\mo\Omega\)의 표현을 낳는다. 기저원소 \(a_i\)들의 표현을 알면 모든 원소의 표현을 알기 때문이다.\textsuperscript{19)}

{\footnotesize\noindent
18a) 더 정확히 말하면, 예전 곱셈법칙을 갖는 새 기호 \(\bar a_i\)를 도입한다. 그러면 \(\bar a_1\Omega+\cdots+\bar a_h\Omega\)는 \(\mo\)와 동형인 부분환을 포함하므로, 나중에 \(\bar a\)를 \(a\)와 동일시할 수 있다.\par
19) 물론 다시 \(P\)-가군준동형인 표현만을 고찰한다. 곧 \(a\mapsto A\)이면 모든 \(\rho\in P\)에 대해 \(a\rho\mapsto A\rho\)이어야 하며, \(\Omega\)에 대해서도 이에 대응하는 조건을 둔다.\par}

아이디얼이나 표현가군, 그리고 이에 대응하는 표현이 대수적으로 닫힌 체로 넘어간 뒤에도 기약으로 남으면 \emph{절대기약}이라 한다.

\(\mo\Omega\)가 근기 없으면 \(\mo\)도 근기 없다. \(\mo\)의 멱영아이디얼 \(\mc\)는 \(\mo\Omega\)의 확장아이디얼 \(\mc\Omega\)를 낳기 때문이다. 그 역은 아래에서 보듯 성립하지 않는다.

\(\mo\Omega\)의 중심은 확대된 중심 \(\mZ\Omega\)와 같다. \(\mo\)가 완전가약이면 \(\mZ\)는 체들의 직접합
\[
        \mZ=\mZ_1+\cdots+\mZ_s
\]
이다. \(\mo\Omega\)로 넘어갈 때 \(\mo\)가 완전가약으로 남으면, 곧 새 멱영아이디얼이 생기지 않으면, 새 중심
\[
        \mZ\Omega=\mZ_1\Omega+\cdots+\mZ_s\Omega
\]
은 \(\Omega\) 안에서 다시 체들의 직접합으로 분해된다. \(\Omega\)가 대수적으로 닫혀 있을 때 이 체들은 § 20의 논평에 따라 \(\Omega\)에 관한 랭크가 1이고 \(\Omega\)와 동형이다.

중심 \(\mZ\)의 표현에는 다음 정리들이 성립한다.

\emph{가환 초복소계 \(\mZ\)의 대수적으로 닫힌 체 \(\Omega\) 안에서의 모든 기약표현은 1차이다. 달리 말하면 기약표현들은 \(\mZ\)에서 \(\Omega\)로 가는 준동형들과 같다.}

\textbf{증명.} \(\mZ\)의 모든 기약표현은 \(\mZ_\Omega\)의 기약표현을 낳고, 따라서 \(\mC\)가 \(\mZ_\Omega\)의 근기일 때 \(\mZ_\Omega/\mC\)의 기약표현도 낳는다. \(\mZ_\Omega/\mC\)는 근기 없는 가환계이므로 § 14 끝부분에 따라 체들의 직접합이다. \(\Omega\)가 대수적으로 닫혀 있으므로 이 체들은 1차이고, 모든 기약표현을 생성한다(§ 20).

동시에 1차 동치표현은 반드시 같으므로(\(\lambda^{-1}\alpha\lambda=\alpha\)) 다음이 따른다. \emph{\(\mZ\)에서 \(\Omega\)로 가는 서로 다른 준동형의 수는 \(\mZ_\Omega/\mC\)의 랭크와 같다.}

마지막 논평은 \(\mZ\)가 \(P\) 위의 체인 특수한 경우에 다음 정리를 준다.

\emph{\(\mZ\)가 가환체이면, \(\mZ\)가 \(P\)의 제1종 확대일 때 그리고 그때에만 \(\mZ\)는 완전가약, 곧 근기 없다.}

실제로 체에 대한 준동형은 동형이다. 그 수는 \(\mZ_\Omega/\mC\)의 랭크와 같으므로, \(\mC=0\)일 때 그리고 그때에만 체차수, 곧 \(\mZ\)가 제1종 확대일 때의 수와 같다.

\(\mZ\)가 완전가약이고
\[
        \mZ=\mZ_1+\cdots+\mZ_t
\]
이면 모든 표현에서 각 체 \(\mZ_\nu\)도 준동형으로 사상된다. 기약표현에서는 이 체들 가운데 하나가 동형으로 사상되고 나머지는 0으로 표현된다(§ 19).

이 정리들을 초복소계에 적용하면 우선 근기 없는 계에 대해 다음을 얻는다.

\emph{\(\mo_\Omega\), 따라서 \(\mo\)가 근기 없는 계이면, \(\mo\)의 모든 절대기약표현에서 중심원소 \(z\)는 대각행렬 \(E\zeta^{(\nu)}\)로 표현되고, \(z\mapsto\zeta^{(\nu)}\)는 \(\mZ\)의 \(\Omega\) 안에서의 1차 표현이다. 이 방식으로 매개되는 \(\mo\)의 절대기약 표현류와 \(\mZ\)의 절대기약 표현류 사이의 관계는 일대일이다. 따라서 이 표현류들의 수는 중심의 랭크와 같다.}\textsuperscript{20)}

{\footnotesize\noindent
20) 대응하는 정리들은 대수적으로 닫힌 체 \(\Omega\) 안에서의 표현뿐 아니라 각각의 자기동형체 안에서의 표현에도 성립하며, 이는 다른 곳에서 전개할 예정이다.\par}

\textbf{증명.} 양쪽 직접분해불능 아이디얼들로의 분해
\[
        \mo_\Omega=\ma_1+\cdots+\ma_t
\]
에는 중심의 체들로의 분해
\[
        \mZ_\Omega=\mZ_1+\cdots+\mZ_t
\]
가 일대일로 대응한다. 여기서 \(\mZ_\Omega\)는 \(\mo_\Omega\)의 중심이다. \(\mo\)의 기약표현에서는 한 \(\ma_\nu\)가 동형으로 사상되고 나머지는 0으로 표현되며, 표현류는 \(\ma_\nu\)에 의해 유일하게 결정된다. \(\ma_\nu\)는 완전행렬환으로 표현되고, 그러한 환의 중심은 대각행렬 \(E\zeta^{(\nu)}\)들로 이루어진다. 대응 \(z\mapsto\zeta^{(\nu)}\)는 \(\mZ\)의 준동형이고, 이때 한 \(\mZ_\nu\)는 동형으로 사상되며 나머지는 0으로 표현된다. 이로써 모두 증명되었다.

근기 없는 \(\mo\)에 언제 역시 근기 없는 \(\mo_\Omega\)가 대응하는가라는 문제에 대해서는 다음도 따른다.

\emph{\(\mZ\)가 \(P\)의 제2종 확대인 성분 \(\mZ_\nu\)를 가지면, \(\mo_\Omega\)는 반드시 근기를 갖는다.}\textsuperscript{21)} \(\mZ_\Omega\)가 이미 근기를 갖기 때문이다.

{\footnotesize\noindent
21) \(\mZ\)가 제1종 성분들만 가지면 \(\mo_\Omega\)는 근기 없게 되며, 이 역시 뒤의 다른 곳에서 증명할 예정이다. 표수가 0인 경우 I. Schur는 이에 동치인 정리, 곧 \(P\) 안에서의 기약표현은 계수영역을 어떻게 확대해도 완전가약으로 남는다는 정리를 이미 증명했다(\emph{Beiträge zur Theorie der Gruppen linearer Substitutionen}, Transact. Am. Math. Soc. 15 (1909), 159쪽).\par}

\subsection*{§ 22. 아벨군에의 적용.}

\(\mG=\mG_1\times\mG_2\times\cdots\times\mG_r\)을 위수가 \(h_i\)인 순환군 \(\mG_i\)들로 분해된 유한 아벨군이라 하자. 그 원소들은
\[
        a=c_1^{\lambda_1}\cdots c_r^{\lambda_r}
\]
이다. \(P\)를 그 표수가 군의 위수 \(h=h_1h_2\cdots h_r\)를 나누지 않는 체라 하자.

군환 \(\mZ\)는 모든 합
\[
        \sum a_\lambda\rho_\lambda
        =
        \sum c_1^{\lambda_1}\cdots c_r^{\lambda_r}\rho_\lambda
        \qquad(\rho_\lambda\in P)
\]
으로 이루어지며, 대응
\[
        \sum \rho_\lambda z_1^{\lambda_1}\cdots z_r^{\lambda_r}
        \longmapsto
        c_1^{\lambda_1}\cdots c_r^{\lambda_r}\rho_\lambda
\]
에 의해 다항식영역 \(P[z_1,\ldots,z_r]\)의 준동형상이다. 따라서 \(\mZ\sim P[z_1,\ldots,z_r]/m\)이고, 쉽게
\[
        m=(z_1^{h_1}-e,\ldots,z_r^{h_r}-e)
\]
를 얻는다.

확대체 \(\Omega\) 안에서 \(z_i^{h_i}-e\)는 서로 다른 일차인자 \(z_i-\epsilon_i\)들로 분해된다. 따라서 \(m\)은 각각 하나의 영점 \((\epsilon_1^{(\nu)},\ldots,\epsilon_r^{(\nu)})\)만 갖는 서로 다른 소아이디얼
\[
        \mP^{(\nu)}
        =(z_1-\epsilon_1^{(\nu)},\ldots,z_r-\epsilon_r^{(\nu)})
\]
들의 곱 또는 교집합이다. 교집합을 취하는 데에는 직접합분해가 대응한다(§ 4).
\[
        \mZ_\Omega=\mZ_1+\cdots+\mZ_h.
\]
여기서 각각 \(\mZ_\nu\simeq\mZ_\Omega/\mP^{(\nu)}\)이고, 따라서 표현
\[
        c_1^{\lambda_1}\cdots c_r^{\lambda_r}
        \longmapsto
        \epsilon_1^{(\nu)\lambda_1}\cdots
        \epsilon_r^{(\nu)\lambda_r},
        \qquad\hbox{또는}\qquad
        a\longmapsto\chi^{(\nu)}(a)
\]
를 매개한다. 이 표현들이 바로 \emph{지표}이다. 그 수는 일반이론이 요구하는 대로 \(h_1h_2\cdots h_r=h\)이다.

\subsection*{§ 23. 초복소계의 행렬식.}

\(\mo=a_1P+\cdots+a_hP\)를 초복소계라 하자. \(P\)에 \(h\)개의 부정원 \(x_1,\ldots,x_h\)를 첨가하여
\[
        \mo^*=a_1P(x)+\cdots+a_hP(x)
\]
를 만든다. \(x_i\)들은 \(a_i\)들과 가환한다고 하며, 이로써 \(\mo^*\)의 계산법칙이 정해진다.

\(\mo^*\) 안에는 \glqq{}\(\mo\)의 일반원소\grqq{}
\[
        w=a_1x_1+\cdots+a_hx_h
\]
가 있다. 한 표현에서 \(a_i\mapsto A_i\)이면 \(w\)에는
\[
        W=A_1x_1+\cdots+A_hx_h
\]
가 대응한다. \(W\)를 그 표현에 속하는 \emph{계행렬}이라 한다. 특히 \(a_i\)들이 군을 이루어 \(\mo\)가 군환이면 \emph{군행렬}이라 한다. 정칙표현의 경우에는 \emph{정칙 계행렬}을 얻는다.

\(W\)의 원소 \(w_{ik}\)들은 \(x\)의 선형형식이다. 따라서 \(n\)차 표현의 경우 \glqq{}\emph{계행렬식}\grqq{} \(|W|\)는 \(n\)차이다. 특히 \emph{정칙 계행렬식}은 \(h\)차이다.

동치표현으로 넘어가도 \(|PWP^{-1}|=|P|\,|W|\,|P^{-1}|=|W|\)이므로 계행렬식은 변하지 않는다. \((a_1,\ldots,a_h)\)에서 새 기저 \((b_1,\ldots,b_h)\)로, \(w=\sum a_ix_i\)에서 \(w=\sum b_iy_i\)로 넘어갈 때에는 예전 원소들에 정칙 치환행렬을 갖는 치환
\[
        x_i=\sum \rho_{ik}y_k
\]
을 시행하여 \(W\)의 새 원소들과 새 행렬식을 얻는다.

표현가군의 합성열이 주어지면 알맞은 기저에서 행렬 \(W\)는
\[
        W=\begin{pmatrix}
        W_1&&0\\
        &W_2&\\
        &\ddots&\\
        W_{ik}&&W_r
        \end{pmatrix}
\]
의 꼴이고,
\[
        |W|=|W_1|\cdot|W_2|\cdots|W_r|.
\]
임의의 표현에 나타나는 \(|W_i|\)들은 \(\mo\), 더 나아가 \(\mc\)가 근기일 때 \(\mo/\mc\)의 정칙표현에 나타나는 것 이외에는 없다.

\emph{\(P\)가 대수적으로 닫혀 있으면, 기약표현에 속하는 행렬식 \(|W_i|\)는 \(x\)들에 관한 기약다항식이고, 서로 동치가 아닌 표현들에는 서로 다른 기약인자들이 속한다.}

\textbf{증명.} \(\mo\)의 모든 기약표현은 \(\mo/\mc\)의 표현이기도 하므로 근기 없는 환 \(\mo/\mc\)에 국한할 수 있다. 그 안에서 행렬단위 \(c_{ik}^{(\nu)}\)들을 기저로 도입하면 일반원소는
\[
        w=\sum c_{ik}^{(\nu)}x_{ik}^{(\nu)}
\]
이다. 기약표현들의 행렬은 \(W_\nu=(x_{ik}^{(\nu)})\)이다. 함수 \(|W_\nu|=|x_{ik}^{(\nu)}|\)들은 잘 알려진 대로 기약이고 서로 다름이 명백하다.

\(|W_\nu|\)들을 계산하려면 \(\mo/\mc\)의 정칙 계행렬을 그 기약인자들로 분해하면 된다. 각 기약인자는 기약표현의 차수가 나타내는 횟수만큼 나타난다. 대응 아이디얼 \(\ml_\nu\)가 합성열에서 그만큼 나타나기 때문이다.

\(\mo\)의 정칙 계행렬을 써도 되지만, 이 경우에는 각 기약인자가 더 자주 나타난다. 곧 대응하는 \(\ml_\nu\)가 좌측아이디얼들의 합성인자로 나타나는 횟수만큼이다. 이 정칙표현에서는 \(\mo\)를 좌측아이디얼로 보았다. \(\mo\)를 우측아이디얼로 보면 두 번째 정칙 계행렬, 곧 프로베니우스의 \glqq{}안티스트로프 행렬\grqq{}이 나온다. 이는 같은 기약인자들, 곧 모든 기약표현의 계행렬식들을 포함하지만 그 지수들은 다를 수 있다(§ 10의 예 참조).

가환계의 계행렬식은 모든 기약표현이 1차이므로 일차인자들로 분해된다. 이 일차인자들은 그 자체가 기약표현이며, 따라서 아벨군의 경우 지표들을 준다. 이 사실이 데데킨트가 비아벨군의 군행렬식을 연구할 때의 출발점이었다.

\subsection*{§ 24. 대각합과 지표.}

초복소계 \(\mo\)의 한 표현 \(\mo\sim\mathfrak D\)에서 원소 \(a\)에 행렬 \(A\)가 대응하면
\[
        \operatorname{Sp}_{\mathfrak D}(a)=\operatorname{Sp}A
\]
로 놓는다.

\emph{대각합은 선형함수이다.}
\[
        \operatorname{Sp}(c+d)
        =\operatorname{Sp}(c)+\operatorname{Sp}(d),
        \qquad
        \operatorname{Sp}(c\alpha)
        =\alpha\operatorname{Sp}(c).
\]

\emph{동치표현들은 같은 대각합을 갖는다.}

\emph{가약표현에서의 대각합은 합성인자들이 매개하는 표현들에서의 대각합의 합이다.}

\textbf{증명.}
\[
        \operatorname{Sp}\begin{pmatrix}
        A_{11}&&0\\
        &A_{22}&\\
        &\ddots&\\
        A_{ik}&&A_{rr}
        \end{pmatrix}
        =\operatorname{Sp}(A_{11})+\operatorname{Sp}(A_{22})
        +\cdots+\operatorname{Sp}(A_{rr}).
\]

\emph{주대각합}은 정칙표현에서의 대각합이다.

\emph{축약대각합}은 서로 다른 기약표현들에서의 대각합의 합이다.

\emph{\(c\)가 극대 멱영아이디얼 \(\mc\)의 원소이면 모든 표현에서 \(\operatorname{Sp}(c)=0\)이다.}

\textbf{증명.} \(\mo/\mc\)의 단순 좌측아이디얼들이 매개하는 표현들만 조사하면 된다. 다른 모든 표현에는 이 합성인자들만 나타나고, 대각합은 가법적으로 합성되기 때문이다.

그런데 \(c\)는 \(\mo/\mc\)의 모든 원소를 소멸시킨다. 따라서 \(\mc\)의 모든 \(c\)에는 영행렬이 대응하고 \(\operatorname{Sp}(c)=0\)이다. 증명 끝.

\(\mo\)가 양쪽단순이고 \(P\)가 대수적으로 닫혀 있어서 \(\mo\)가 행렬환
\[
        \mo=\sum c_{ik}P
\]
이라고 하자. \(\mo\)의 기약표현에서 원소 \(a=\sum c_{ik}\alpha_{ik}\)에는 행렬 \((\alpha_{ik})\)가 대응한다. 따라서
\[
\begin{aligned}
        \hbox{축약대각합 }\operatorname{Sp}_{\ml}(a)
          &=\sum_i\alpha_{ii},\\
        \hbox{주대각합 }\operatorname{Sp}_{\mo}(a)
          &=n\cdot\hbox{축약대각합}
           =n\sum_i\alpha_{ii}.
\end{aligned}
\]
특히
\[
\begin{aligned}
        \operatorname{Sp}(c_{ik})&=0\quad(i\ne k),\\
        \operatorname{Sp}_{\ml}(c_{ii})&=e,\\
        \operatorname{Sp}_{\mo}(c_{ii})&=n e.
\end{aligned}
\]

\emph{\(P\)에서 대수적으로 닫힌 체 \(\Omega\)로 넘어가도 주대각합은 변하지 않는다.} 같은 기저에서 계산할 수 있기 때문이다. 축약대각합에는 이 명제가 성립하지 않는다.

근기 없는 계의 원소 \(a\)가 절대기약표현들에서 갖는 대각합을 \emph{지표}라 하고 \(\chi(a)\), 또는 어느 표현을 뜻하는지 나타내려면 \(\chi^{(\nu)}(a)\)로 쓴다.

\(n_\nu\)차 기약표현에서 중심원소들은 § 21에 따라 대각행렬 \(E\zeta^{(\nu)}\)로 표현된다. 여기서 \(z\mapsto\zeta^{(\nu)}\)는 중심의 1차 표현, 곧 중심에서 \(\Omega\)로 가는 준동형이다. 따라서 그 대각합은 \(n_\nu\zeta^{(\nu)}\)이다. \emph{그러므로 중심의 준동형 \(\Theta\)와 지표 \(\chi\)는 관계}
\srcnumdisplay{(1)}{%
        \chi(z)=n_\nu\Theta(z)}
\emph{로 연결된다.} 가환인 경우에는 \(n_\nu=1\)이고 지표 자체가 준동형을 준다(§ 22 참조). 이하에서 언제나 가정하듯 체 \(\Omega\)의 표수가 0이면 (1)을 \(n_\nu\)로 나눌 수 있다.
\[
        \Theta(z)=\frac{\chi(z)}{n_\nu}.
\]
\(\Theta(z)\)의 준동형 성질은 식
\[
        \frac{\chi(z)}{n_\nu}
        +\frac{\chi(z')}{n_\nu}
        =\frac{\chi(z+z')}{n_\nu},
\]
\[
        \frac{\chi(z)}{n_\nu}
        \cdot\frac{\chi(z')}{n_\nu}
        =\frac{\chi(zz')}{n_\nu}
\]
으로 표현된다.

\emph{한 표현류는 행렬들의 대각합만으로 이미 유일하게 결정된다.} (모든 행렬의 대각합을 알려면 주어진 표현에서 기저원소들의 대각합만 알면 충분함은 물론이다.)

\textbf{증명.} 표현가군 \(\mR\)은 각 기약 표현가군 \(\mM_\nu\)가 직접가수로 몇 번 나타나는지를 알면 결정된다. 그 횟수가 \(p_\nu\)이면 표현에서 \(e^{(\nu)}\)의 대각합은 \(p_\nu n_\nu\)이다. 여기서 \(n_\nu\)는 기약표현의 차수이다. 따라서
\[
        p_\nu=\frac{\operatorname{Sp}e^{(\nu)}}{n_\nu}.
\]

\subsection*{§ 25. 판별식.}

\(\mo=a_1P+\cdots+a_hP\)를 초복소계라 하자.

\noindent\emph{판별행렬}은 그 원소들이 주대각합 \(\operatorname{Sp}(a_i a_k)\)인 행렬이다.

\noindent\emph{축약판별행렬}은 대수적으로 닫힌 체에서 만든 축약대각합에 관하여 같은 방식으로 얻는 행렬이다.

이 행렬의 행렬식을 \emph{판별식}, 각각 \emph{축약판별식}이라 한다.

\noindent\emph{바탕체를 확대해도 판별식은 변하지 않는다.}

\noindent\emph{다른 기저로 넘어가면 판별식에는 변환행렬의 행렬식의 제곱이 곱해진다.}

\textbf{증명.}
\begin{align*}
(a_1,\ldots,a_h)&=(b_1,\ldots,b_h)P,\\
(a_i a_1,\ldots,a_i a_h)&=(a_i b_1,\ldots,a_i b_h)P,\\
\bigl(\operatorname{Sp}(a_i a_1),\ldots,
      \operatorname{Sp}(a_i a_h)\bigr)
&=\bigl(\operatorname{Sp}(a_i b_1),\ldots,
        \operatorname{Sp}(a_i b_h)\bigr)P.
\end{align*}
따라서 행렬로 쓰면
\srcnumdisplay{(1)}{%
\bigl(\operatorname{Sp}(a_i a_k)\bigr)
=\bigl(\operatorname{Sp}(a_i b_k)\bigr)P.}
마찬가지로 \(\widetilde P\)가 전치행렬이면
\[
\begin{pmatrix}a_1\\[-1mm]\vdots\\[-1mm]a_h\end{pmatrix}
=\widetilde P
\begin{pmatrix}b_1\\[-1mm]\vdots\\[-1mm]b_h\end{pmatrix}
\]
이고, 앞에서와 같이
\srcnumdisplay{(1a)}{%
\bigl(\operatorname{Sp}(a_i b_k)\bigr)
=\widetilde P\bigl(\operatorname{Sp}(b_i b_k)\bigr).}
(1)과 (1a)에서
\[
\bigl(\operatorname{Sp}(a_i a_k)\bigr)
=\widetilde P\cdot
 \bigl(\operatorname{Sp}(b_i b_k)\bigr)\cdot P
\]
를 얻는다. 행렬식을 취하면 주장이 따른다.

따라서 판별식은 \(P\)의 제곱수 인자를 제외하면 정해질 뿐이다. 그러나 판별식이 0인지 아닌지는 불변성질이다.

(1)에서는 다음도 따른다. \emph{서로 다른 두 기저를 써서 \(\left|\bigl(\operatorname{Sp}(a_i b_k)\bigr)\right|\)로 판별식을 0이 아닌 한 인자를 제외하고 구할 수도 있다.}

\noindent\emph{\(\mo\)가 멱영아이디얼 \(\mc\)를 가지면 판별식은 0이다.}

\textbf{증명.} \(\mo\)의 기저로 \((c_1,\ldots,c_t,d_{t+1},\ldots,d_h)\)를 택하되, \((c_1,\ldots,c_t)\)가 \(\mc\)의 기저가 되게 하자. 판별식은
\[
\left|\begin{array}{cc}
\operatorname{Sp}(c_i c_k)&\operatorname{Sp}(c_i d_k)\\
\operatorname{Sp}(d_i c_k)&\operatorname{Sp}(d_i d_k)
\end{array}\right|
=
\left|\begin{array}{cc}
0&0\\
0&\operatorname{Sp}(d_i d_k)
\end{array}\right|=0.
\]
(축약판별식에도 이 명제가 성립한다.)

\textbf{따름정리.} \emph{대수적으로 닫힌 체로 넘어간 뒤 멱영아이디얼이 나타나는 경우에도 판별식은 0이다.}

\(\mo=\ma_1+\cdots+\ma_s\)가 양쪽 직접합이고, \(M,M_1,\ldots,M_s\)가 각각 환 \(\mo,\ma_1,\ldots,\ma_s\)의 판별행렬이라 하자. 그러면
\[
M=\begin{pmatrix}
M_1&0& &0\\
 &M_2& &\\
 & &\ddots&\\
0& & &M_s
\end{pmatrix},
\]
따라서 \(|M|=|M_1|\,|M_2|\cdots|M_s|\)이다.

\textbf{증명.} \(\ma_1,\ldots,\ma_s\)의 기저들을 합쳐 \(\mo\)의 기저를 택하자. \(a\in\ma_i\)이면
\[
        \operatorname{Sp}_{\mo}(a)
        =\operatorname{Sp}_{\ma_i}(a)
\]
이다. 양쪽 모두 각각 환 \(\mo\), \(\ma_i\)에서의 주대각합을 뜻한다. 또한 \(a_i\in\ma_i\), \(a_k\in\ma_k\)이고 \(i\ne k\)이면 \(a_i a_k=0\), 따라서 \(\operatorname{Sp}(a_i a_k)=0\)이다. 여기서 주장이 따르며, 축약판별식에 대해서도 똑같이 따른다.

행렬환 \(\sum c_{ik}P\)의 \emph{판별식}을 만들기 위해 두 서로 다른 기저를 택하자. 한 번은 \(c_{ik}\)들을 첫째 지수에 따라, 다른 한 번은 둘째 지수에 따라 배열한다. 곱셈표는 다음과 같다.
\[
\begin{array}{c|c|c|c}
& \overbrace{c_{11}\;\cdots\;c_{1n}}^{r_1}
& \overbrace{c_{21}\;\cdots\;c_{2n}}^{r_2}
& \cdots \\ \hline
\left.\begin{array}{c}c_{11}\\[-1mm]\vdots\\[-1mm]c_{n1}\end{array}\right\}l_1
& (c_{ik})&0&0\\[1.3em]\hline
\left.\begin{array}{c}c_{12}\\[-1mm]\vdots\\[-1mm]c_{n2}\end{array}\right\}l_2
&0&(c_{ik})&0\\[1.3em]\hline
\vdots&0&0&(c_{ik})
\end{array}
\]
판별식은
\[
D=\left|\begin{array}{cccc}
\operatorname{Sp}(c_{11})&0&\cdots&0\\
0&\operatorname{Sp}(c_{22})&\cdots&0\\
\vdots&\vdots&\ddots&\vdots\\
0&0&\cdots&\operatorname{Sp}(c_{nn})
\end{array}\right|^{n}
=\begin{cases}
 e,&\hbox{축약대각합의 경우},\\
 n^{n^2}e,&\hbox{주대각합의 경우}.
\end{cases}
\]

\textbf{따름정리.} \emph{\(\mo\)가 대수적으로 닫힌 확대체에서 차수가 \(n_1,\ldots,n_s\)인 행렬환들로 분해되면 판별식은}
\[
        D=n_1^{n_1^2}n_2^{n_2^2}\cdots n_s^{n_s^2}e
\]
\emph{이고, 축약판별식은}
\[
        D_{\rm red}=e
\]
\emph{이다.}

따라서 근기 없는 계에서는 축약판별식이 언제나 0이 아니다. 다른 판별식은 \(n_i\) 가운데 어느 것도 체의 표수로 나누어지지 않을 때에만 0이 아니다. 그러므로 다음이 성립한다.

\noindent\emph{축약판별식이 0이 아니라는 것은 체 \(P\)를 대수적으로 닫은 뒤 근기가 없는 계가 되기 위한 필요충분조건이다.}

\noindent\emph{표수가 0인 경우 판별식이 0이 아니라는 것은 체 \(P\)를 대수적으로 닫은 뒤 근기가 나타나지 않기 위한 필요충분조건이다.}\textsuperscript{22)}

{\footnotesize\noindent
22) 가환계에 대해서는 E. Noether, \emph{Diskriminantensatz für Ordnungen\ldots}, J. f. M. 157 (1927), 82--104쪽, §§ 4--6 참조. 그곳의 증명법은 같지만 세부적으로 더 복잡하다. 한 기저에서 다른 기저로 넘어가는 과정(§ 4, 4)은 여기, 곧 이 절 첫머리에서 준 과정으로 바꾸어야 한다. 행렬식 \(\left|(e_i e_k)\right|\)는 0이기 때문이다.\par}

\subsection*{§ 26. 군환의 분류.}

\noindent\emph{\(a_1,\ldots,a_h\)가 유한군의 원소들이고 표수가 \(h\)의 약수가 아닌 체에서 군환을 만들면, 판별식 \(D\ne0\)이고 따라서 군환은 근기 없는 환이다.}

\textbf{증명.} 먼저 다음을 보이자. 이제부터 \(\operatorname{Sp}\)는 주대각합을 뜻한다.
\[
        \operatorname{Sp}(e)=he,
        \qquad
        \operatorname{Sp}(a_i)=0\quad(a_i\ne e).
\]
정칙표현에서
\[
e\longmapsto
\begin{pmatrix}
e&0&\cdots&0\\
0&e&\cdots&0\\
\vdots&\vdots&\ddots&\vdots\\
0&0&\cdots&e
\end{pmatrix},
\qquad
\operatorname{Sp}(e)=he
\]
이다. \(a_i\ne e\)이면
\[
        a_i a_k=a_\lambda,
        \qquad \lambda\ne k.
\]
따라서 곱 \(a_i a_1,\ldots,a_i a_h\)를 \(a_1,\ldots,a_h\)로 나타내는 행렬은 주대각선의 모든 원소가 0이고, 그러므로 \(\operatorname{Sp}(a_i)=0\)이다.

이제 두 기저 \(a_1,\ldots,a_h\)와 \(a_1^{-1},\ldots,a_h^{-1}\)를 택한다. 행렬 \(\bigl(\operatorname{Sp}(a_i a_k^{-1})\bigr)\)은
\[
\begin{pmatrix}
he&0&\cdots&0\\
0&he&\cdots&0\\
\vdots&\vdots&\ddots&\vdots\\
0&0&\cdots&he
\end{pmatrix}
\]
이다. 따라서
\[
        D=h^h e\ne0.
\]
증명 끝.

군원소의 \emph{류}란 군환에서 만든 합
\srcnumdisplay{(1)}{K_a=\sum_s s^{-1}as}
을 말한다. 여기서는 서로 다른 \(s^{-1}as\)들에 대해서만 합한다. \emph{류 \(K_a\)들은 중심원소이다.} 모든 군원소와 가환하기 때문이다. \emph{류 \(K_a\)들은 중심을 생성한다.} 실제로 군환의 원소 \(\sum_i a_i\varrho_i\)가 모든 \(s\)와 가환하면
\[
\begin{aligned}
\sum_i a_i\varrho_i
&=s^{-1}\biggl(\sum_i a_i\varrho_i\biggr)s
 =\sum_i(s^{-1}a_i s)\varrho_i\\
&=\frac1h\sum_i\biggl(\sum_s s^{-1}a_i s\biggr)\varrho_i
 =\frac1h\sum_i\biggl(\frac h{h_i}K_i\biggr)\varrho_i.
\end{aligned}
\]
여기서 \(h\)는 군원소의 수이고 \(h_i\)는 류 \(K_i\)의 원소 수이다.

따라서 중심의 랭크는 켤레군원소류의 수와 같고, 그러므로 \emph{절대기약표현의 수 역시 이 류의 수와 같다.}

행렬 \(A\)와 \(S^{-1}AS\)는 같은 대각합을 갖는다. 따라서 군원소 \(a\)와 \(s^{-1}as\)는 같은 대각합을 갖는다.

이제 (1)의 양변에서 대각합을 취하면
\[
        \chi(K_i)=h_i\chi(a_i)
\]
를 얻는다. 말로 쓰면 다음과 같다.

\noindent\emph{군원소의 지표는 그 류의 지표를 그 류의 원소 수로 나눈 것과 같다.}

\begin{center}
(1928년 8월 12일 접수.)
\end{center}

\emph{\(\mG\)가 완전가약이면 모든 정규부분군은 직접인자이고, 다른 인자는 주어진 \(\mG\)의 곱분해에 나타나는 단순군들의 곱으로 택할 수 있다. 정규부분군 \(\mH\)도 완전가약이다.}

\emph{증명.} 
\begin{equation}
        \mG=\mH\mA_1\mA_2\cdots\mA_n. \tag{3}
\end{equation}
\(\mH_i=\mH\mA_1\cdots\mA_i\)로 놓으면 \(\mH_{i+1}=\mH_i\mA_{i+1}\)이다. 이제 \(\mA_{i+1}\le\mH_i\), 따라서 \(\mH_{i+1}=\mH_i\)이면 (3)에서 인자 \(\mA_{i+1}\)을 생략한다. 그렇지 않으면 \(\mH_i\cap\mA_{i+1}\)은 \(\mA_{i+1}\)의 진정규부분군이므로 \(\mE\)와 같고, \(\mH_i\times\mA_{i+1}\)은 직접곱이다. 불필요한 인자들을 생략하면 (3)은
\[
        \mG=\mH\times\mA_{i_1}\times\cdots\times\mA_{i_r}
\]
이 되어 \(\mH\)가 직접인자임을 보인다. 또한
\[
        \mH\simeq
        \mG/\mA_{i_1}\times\cdots\times\mA_{i_r}
        \simeq
        \mA_{j_1}\times\cdots\times\mA_{j_\mu},
\]
여기서 \(\mA_{j_1},\ldots,\mA_{j_\mu}\)는 나머지 \(\mA_i\)들이다. 따라서 \(\mH\)도 완전가약이다.

\emph{따름정리.} \emph{완전가약군을 직접분해불능 인자들로 나타내는 모든 분해는 단순인자들로의 분해이며, 따라서 합성열을 낳는다. 이로부터 인자들이 동형을 제외하면 유일하게 결정된다는 것이 따른다.}
